% Section 4, continued: the proof of Remark 4.14 (Seidel's seven graphs).
% Route: cgss/prerequisites.md, Section 7.3 (decision 3 of Section 13).
% Macros of this file, in case macros.tex does not provide them.
% \Ecal_t: the edge set E_t of the graph \Delta_t on {1,...,8} (t = 6, 7, 8).
\providecommand{\Ecal}{\mathcal E}
% The 4-clique profile of a graph, as a multiset.
\providecommand{\cprof}{\opk{prof}_4}
% The symmetric difference of two sets.
\providecommand{\sd}{\mathbin{\triangle}}

\subsection{The proof of Remark 4.14}\label{sec:rem414}
\begingroup
\emergencystretch=3em

We will prove Remark 4.14 by the route of \path{cgss/prerequisites.md}, Section 7.3, which
follows the computer-free part of the proof of [BCN] Theorem 3.12.2 in the language of line
systems, and replaces the enumeration of [BCN] Proposition 3.11.2 by three switching arguments.
A graph $X$ as in the remark is represented by a family $(\Zx_u)$ in $\typ E_8$. We will first
pin down its parameters (Computation C18), and then study the line system $R$ of the vectors of
squared norm 2 in the $\ints$-span of the $\Zx_u$. It is of type $\typ E_6$, $\typ E_7$ or
$\typ E_8$, and the vector $\Zw=\frac1{k+2}\sum_u\Zx_u$ has integral inner products with it. The
simple roots of Computation C21 then exclude two parameter sets. In the three cases where
$(\Zw,\Zw)=2$, $\Zw$ is a root, and $X$ is a switching of $T(5)$, of $T(6)$ with an isolated
vertex, or of $T(8)$, by the edge set of a regular graph; Computation C22 names the result. The
two remaining parameter sets are those of the Clebsch and the Schl\"afli graphs, which are
unique. We will write $n$, $k$, $a$, $c$ for the parameters of a strongly regular graph, as in
the rest of the blueprint. A result tagged \new\ departs from Section 7.3 of
\path{cgss/prerequisites.md}, or from the prerequisite that it cites, and the departure is
stated after the tag; the untagged results follow it.

\begingroup\raggedright
\medskip\noindent\textbf{Facts from Mathlib and the project.}\label{facts:414}
We will use the following declarations of Mathlib (commit \path{082e2d37}):
\path{SimpleGraph.IsSRGWith} with its fields,
\path{SimpleGraph.IsSRGWith.matrix_eq} ($A^2=kI+aA+c\,A(\comp X)$),
\path{SimpleGraph.one_add_adjMatrix_add_compl_adjMatrix_eq_of_one} with
\path{SimpleGraph.compl_adjMatrix_eq_adjMatrix_compl} ($I+A+A(\comp X)=J$),
\path{SimpleGraph.IsSRGWith.param_eq} ($k(k-a-1)=c(n-k-1)$ for $0<n$, in $\mathbb N$ with truncated
subtraction),
\path{SimpleGraph.IsSRGWith.compl} (the complement has parameters
$(n,n-k-1,n-2k+c-2,n-2k+a)$),
\path{SimpleGraph.IsSRGWith.card_neighborFinset_union_of_not_adj} ($|N(u)\cup N(w)|=2k-c$ for
distinct non-adjacent $u,w$),
\path{SimpleGraph.IsRegularOfDegree.compl},
\path{SimpleGraph.sum_degrees_eq_twice_card_edges}, \path{SimpleGraph.IsNClique.map} (an
embedding maps cliques to cliques), \path{compl_compl} ($\comp{\comp X}=X$),
\path{Matrix.exists_mulVec_eq_zero_iff} (a square matrix is singular exactly when it has a
nonzero null vector), \path{Matrix.inv_def} and \path{Matrix.mul_adjugate},
\path{Matrix.det_gram_ne_zero_iff_linearIndependent},
\path{LinearMap.IsProj.trace} with \path{LinearMap.isProj_range_iff_isIdempotentElem} (an
idempotent has rank equal to its trace), \path{abs_real_inner_le_norm} (the Cauchy--Schwarz
inequality), \path{exists_linearIndependent} (a spanning set contains a basis) and
\path{Submodule.finrank_le}. From the project we use
\path{SimpleGraph.leastEigenvalue_eq_neg_two_iff} of
\path{LineGraphs/Spectral/Positivity.lean} ($\lambda_{\min}(X)=-2$ exactly when $2I+A$ is
positive semidefinite and singular), the checkers \path{Computations.checkGraphIso}, \path{Computations.cliqueProfile} and
\path{Computations.cliqueCount} with their soundness theorems
\path{Computations.nonempty_graphOfMatrix_iso}, \path{Computations.checkNonisomorphic_sound} and
\path{Computations.cliqueCount_sound},
and the invariance \path{SimpleGraph.Iso.cliqueProfile_eq}, in
\path{LineGraphs/Computations/CliqueCertificates.lean}. Finally, from the Lean of Computation
C4 we use the transitivity on ordered stars: for $t\in\{6,7,8\}$ and roots $\Za,\Zb$ of
$\typ E_t$ with $(\Za,\Zb)=-1$, there is a linear isometry $g$ of $\re^8$ with
$g(\Phi(\typ E_t))=\Phi(\typ E_t)$ that carries the fixed star $(\Za_0,\Zb_0)$ of Lean to
$(\Za,\Zb)$; this is \leancite{LineSystem.E₈.exists_isometry_star}{\texttt{LineSystem.E\ensuremath{{}_8}.exists\_isometry\_star}}
and its analogues for $\typ E_7$ and $\typ E_6$, in
\path{LineGraphs/Computations/C4E8Stars.lean} and its analogues, whose trees of reflections are
checked by \texttt{decide +kernel}. It replaces LN5 of \path{cgss/prerequisites.md}.
\par
\endgroup

\medskip
Throughout, $\Zc=\frac12\sum_{i=1}^8\Ze_i$ is the root of model (ii) of $\typ E_8$ fixed before
Lemma~\ref{lem:4-switch}, and for a 2-subset $P=\{i,j\}$ of $\{1,\dots,8\}$ we write
$\Ze_P=\Ze_i+\Ze_j$. We will write $\Ecal_6$ for the set of 2-subsets of $\{1,\dots,5\}$,
$\Ecal_7$ for the set of 2-subsets of $\{1,\dots,6\}$ together with $\varepsilon=\{7,8\}$, and
$\Ecal_8$ for the set of all 2-subsets of $\{1,\dots,8\}$; $\Delta_t$ is the graph on
$\{1,\dots,8\}$ with edge set $\Ecal_t$. Thus $L(\Delta_6)\cong T(5)$, $L(\Delta_7)$ is $T(6)$ together
with the isolated vertex $\varepsilon$, and $L(\Delta_8)=T(8)$. For a graph $Z$ and a vertex $i$ we
write $\deg_Z(i)$ for its degree, and for a set $H$ of 2-subsets we write $\deg_H(i)$ for the
number of members of $H$ that contain $i$.

\begin{namedresult}{Definition}{(switching)}\label{def:414-switch}
 \lean{SimpleGraph.switch}
 \uses{}
 \assume{$X$ is a graph and $U\sbs V(X)$.}
 \pflet{The \emph{switched graph} $X^U$ has vertex set $V(X)$, and distinct $u,w$ are adjacent
 in $X^U$ when either they are adjacent in $X$ and both or neither lie in $U$, or they are not
 adjacent in $X$ and exactly one of them lies in $U$. For a graph $Y$ and $H\sbs E(Y)$ we
 write $L(Y)^H$ for the switching of the line graph $L(Y)$ with respect to $H$.}
 This is SW1 of \path{cgss/prerequisites.md}, with $U$ in place of its $Y$, since $Y$ is a
 graph here.
\end{namedresult}

\begin{namedresult}{Lemma}{(switching)}\label{lem:414-switch}
 \lean{SimpleGraph.switch_compl}\lean{SimpleGraph.Iso.switch}\lean{SimpleGraph.Iso.lineGraphSwitch}
 \uses{def:414-switch}
 \assume{$X$ and $X'$ are graphs, $U\sbs V(X)$, $Y$ and $Y'$ are graphs and $H\sbs E(Y)$.}
 \prove{\begin{pfenum}
  \item $X^U=X^{V(X)\setminus U}$.
  \item If $\sigma$ is an isomorphism from $X$ to $X'$, then $\sigma$ is an isomorphism from
   $X^U$ to $X'^{\sigma(U)}$.
  \item If $\sigma$ is an isomorphism from $Y$ to $Y'$, then $e\mapsto\sigma(e)$ is an
   isomorphism from $L(Y)^H$ to $L(Y')^{\sigma(H)}$.
 \end{pfenum}}
\end{namedresult}
\begin{proof}
 \pfsketch\ We will read each assertion off Definition~\ref{def:414-switch}.
 \step{s414swcompl}{Assertion 1 holds.}
 \begin{proof}\pf\ A pair has exactly one member in $U$ if and only if it has exactly one member
  in $V(X)\setminus U$; apply Definition~\ref{def:414-switch}.~\qed\end{proof}
 \step{s414swiso}{Assertion 2 holds.}
 \begin{proof}\pf\ For distinct $u,w$, $\sigma$ preserves adjacency, and $u\in U$ exactly when
  $\sigma(u)\in\sigma(U)$; so the condition of Definition~\ref{def:414-switch} for $u,w$ in $X$
  and $U$ is that for $\sigma(u),\sigma(w)$ in $X'$ and $\sigma(U)$.~\qed\end{proof}
 \step{s414swline}{$e\mapsto\sigma(e)$ is an isomorphism from $L(Y)$ to $L(Y')$.}
 \begin{proof}\pf\ It is a bijection from $E(Y)$ onto $E(Y')$, since $\sigma$ is an isomorphism,
  and $|\sigma(e)\cap\sigma(f)|=|e\cap f|$ since $\sigma$ is injective; two edges are adjacent in
  the line graph when they are distinct and meet.~\qed\end{proof}
 \qedstep
 \begin{proof}\pf\ Assertions 1 and 2 are \stepref{s414swcompl} and \stepref{s414swiso};
  assertion 3 is assertion 2 applied to the isomorphism of \stepref{s414swline} and the set
  $H$.~\qed\end{proof}
\end{proof}

\begin{namedresult}{Lemma}{(isomorphisms and complements)}\label{lem:414-iso}
 \lean{SimpleGraph.Iso.compl}\lean{SimpleGraph.IsSRGWith.of_iso}
 \uses{}
 \assume{$X$ and $X'$ are graphs, and $\sigma$ is an isomorphism from $X$ to $X'$.}
 \prove{\begin{pfenum}
  \item $\sigma$ is an isomorphism from $\comp X$ to $\comp{X'}$.
  \item If $X$ is strongly regular with parameters $(n,k,a,c)$, then so is $X'$.
 \end{pfenum}}
 Mathlib has the complement of an embedding, \path{SimpleGraph.Embedding.complEquiv}, but neither
 assertion.
\end{namedresult}
\begin{proof}
 \pfsketch\ An isomorphism is a bijection that preserves adjacency and non-adjacency.
 \step{s414isocompl}{Assertion 1 holds.}
 \begin{proof}\pf\ $\sigma$ is a bijection, and for distinct $u,w$ the vertices $\sigma(u)$,
  $\sigma(w)$ are distinct and adjacent in $X'$ exactly when $u,w$ are adjacent in $X$; distinct
  vertices are adjacent in the complement exactly when they are not adjacent in the
  graph.~\qed\end{proof}
 \step{s414isosrg}{Assertion 2 holds.}
 \begin{proof}\pf\ $|V(X')|=|V(X)|$, $\sigma$ maps the neighbours of $u$ onto those of $\sigma(u)$,
  and the common neighbours of $u,w$ onto those of $\sigma(u),\sigma(w)$.~\qed\end{proof}
 \qedstep
 \begin{proof}\pf\ \stepref{s414isocompl} and \stepref{s414isosrg}.~\qed\end{proof}
\end{proof}

\subsubsection*{Parameters}

\begin{namedresult}{Lemma}{(cocktail party parameters)}\label{lem:414-CP}
 \lean{SimpleGraph.IsRegularOfDegree.nonempty_iso_cocktailParty}
 \uses{def:CP,facts:414}
 \assume{$m\ge1$, and $X$ is a $(2m-2)$-regular graph on $2m$ vertices.}
 \prove{$X\cong CP(m)$.}
\end{namedresult}
\begin{proof}
 \pfsketch\ The complement of $X$ is a perfect matching.
 \step{s414cpcompl}{The complement $\comp X$ is 1-regular.}
 \begin{proof}\pf\ \path{SimpleGraph.IsRegularOfDegree.compl}: its degree is
  $2m-1-(2m-2)=1$.~\qed\end{proof}
 \step{s414cpmatch}{\pick\ vertices $p_1,q_1,\dots,p_m,q_m$, which are all the vertices of $X$,
  such that the edges of $\comp X$ are the pairs $\{p_i,q_i\}$.}
 \begin{proof}\pf\ By \stepref{s414cpcompl} each vertex lies in exactly one edge of $\comp X$, so
  the edges of $\comp X$ partition $V(X)$ into pairs; there are $2m/2=m$ of
  them.~\qed\end{proof}
 \qedstep
 \begin{proof}\pf\ The map $(i,0)\mapsto p_i$, $(i,1)\mapsto q_i$ is a bijection by
  \stepref{s414cpmatch}. Distinct vertices are adjacent in $X$ exactly when they are not an edge
  of $\comp X$, that is, when their indices differ; this is adjacency in $CP(m)$
  (Definition~\ref{def:CP}).~\qed\end{proof}
\end{proof}

\begin{namedresult}{Lemma}{(strongly regular graphs with eigenvalue $-2$)}\label{lem:414-srg}
 \lean{SimpleGraph.IsSRGWith.params_of_det_two_add_adjMatrix_eq_zero}
 \lean{SimpleGraph.IsSRGWith.rank_two_add_adjMatrix_of_det_eq_zero}
 \uses{facts:414}
 \assume{\begin{pfenum}
  \item $X$ is connected and strongly regular with parameters $(n,k,a,c)$.
  \item $\det(2I+A)=0$, where $A=A(X)$.
 \end{pfenum}}
 \prove{\begin{pfenum}
  \item $1\le k\le n-2$, $a\le k-1$ and $1\le c\le k$.
  \item $2k-c\le n-2$.
  \item $k-2a+c=4$.
  \item $(a-c+4)(\rk(2I+A)-1)=2n-k-2$.
 \end{pfenum}}
 \new\ This is SR4 of \path{cgss/prerequisites.md}, in a form without the spectral theorem:
 the multiplicity $m_1$ of the second eigenvalue is replaced by $\rk(2I+A)-1$, through a
 spectral idempotent. Assertion 2 is the counting condition that the search of C18 needs
 (\path{cgss/certificates/c18_parameter_search.sage}).
\end{namedresult}
\begin{proof}
 \pfsketch\ We will derive the linear relation from a null vector of $2I+A$, and then find an
 idempotent multiple of $2I+A-\frac{k+2}nJ$.
 \pflet{$Q=2I+A$, and let $J$ and $\Zj$ be the all-ones matrix and vector.}
 \step{s414srgnc}{$k\ge1$, and $X$ is not complete.}
 \begin{proof}\pf\ If $k=0$ then $Q=2I$, and if $X$ is complete then $Q=I+J$, which is positive
  definite since $(\Zz,(I+J)\Zz)=|\Zz|^2+(\Zz,\Zj)^2$; in both cases $\det Q\ne0$, against
  \asm[2]{0}.~\qed\end{proof}
 \step{s414srgpairs}{There are an edge $uw$ and distinct non-adjacent vertices $u',w'$.}
 \begin{proof}\pf\ $k\ge1$ and $X$ is not complete, by \stepref{s414srgnc}.~\qed\end{proof}
 \step{s414srgak}{$a\le k-1$.}
 \begin{proof}\pf\ The $a$ common neighbours of the edge $uw$ of \stepref{s414srgpairs} lie in
  $N(u)\setminus\{w\}$, which has $k-1$ elements.~\qed\end{proof}
 \step{s414srgck}{$c\le k$.}
 \begin{proof}\pf\ The $c$ common neighbours of $u',w'$ of \stepref{s414srgpairs} lie in $N(u')$,
  which has $k$ elements.~\qed\end{proof}
 \step{s414srgunion}{$2k-c\le n-2$.}
 \begin{proof}\pf\ $N(u')\cup N(w')$, for $u',w'$ of \stepref{s414srgpairs}, has $2k-c$ elements by
  \path{SimpleGraph.IsSRGWith.card_neighborFinset_union_of_not_adj}, and lies in
  $V(X)\setminus\{u',w'\}$, since $u'\not\sim w'$.~\qed\end{proof}
 \step{s414srgkn}{$k\le n-2$.}
 \begin{proof}\pf\ $k\le2k-c$ by \stepref{s414srgck}; apply \stepref{s414srgunion}.~\qed\end{proof}
 \step{s414srgc}{$c\ge1$.}
 \begin{proof}\pf\ Take non-adjacent distinct $u',w'$ (\stepref{s414srgnc}) and a shortest path
  $u'=u_0,u_1,\dots,u_d=w'$ in the connected graph $X$, so $d\ge2$. Then $u_0$ and $u_2$ are
  distinct and non-adjacent, by minimality, and $u_1$ is a common
  neighbour.~\qed\end{proof}
 \step{s414srgcompl}{$A(\comp X)=J-I-A$.}
 \begin{proof}\raggedright\pf\ \path{SimpleGraph.one_add_adjMatrix_add_compl_adjMatrix_eq_of_one} and
  \path{SimpleGraph.compl_adjMatrix_eq_adjMatrix_compl}.~\qed\end{proof}
 \step{s414srgmat}{$A^2=kI+aA+c(J-I-A)$, $AJ=JA=kJ$ and $J^2=nJ$.}
 \begin{proof}\pf\ \path{SimpleGraph.IsSRGWith.matrix_eq} gives $A^2=kI+aA+c\,A(\comp X)$; apply
  \stepref{s414srgcompl}. $AJ=JA=kJ$ by regularity, and $J^2=nJ$ since each entry of $J^2$ is a
  sum of $n$ ones.~\qed\end{proof}
 \step{s414srgnull}{\pick\ $\Zz\ne0$ with $Q\Zz=0$; then $A\Zz=-2\Zz$ and $J\Zz=0$.}
 \begin{proof}\pf\ \path{Matrix.exists_mulVec_eq_zero_iff} and \asm[2]{0} give $\Zz$. Then
  $k(\Zj,\Zz)=(A\Zj,\Zz)=(\Zj,A\Zz)=-2(\Zj,\Zz)$ by \stepref{s414srgmat} and the symmetry of $A$,
  so $(\Zj,\Zz)=0$, since $k+2>0$; and $J\Zz=(\Zj,\Zz)\Zj$.~\qed\end{proof}
 \step{s414srglin}{$k-2a+c=4$.}
 \begin{proof}\pf\ Apply both sides of \stepref{s414srgmat} to $\Zz$ of \stepref{s414srgnull}:
  $4\Zz=k\Zz-2a\Zz+c(0-\Zz+2\Zz)=(k-2a+c)\Zz$, and $\Zz\ne0$.~\qed\end{proof}
 \pflet{$\eta=a-c+2$.}
 \step{s414srgtheta}{$k-c=2\eta$, $\eta\ge0$, and $(A-\eta I)Q=cJ$.}
 \begin{proof}\pf\ The first is \stepref{s414srglin}, and $\eta\ge0$ follows with $c\le k$
  (\stepref{s414srgck}). By \stepref{s414srgmat},
  $(A-\eta I)(A+2I)=A^2+(2-\eta)A-2\eta I=(k-c-2\eta)I+(a-c+2-\eta)A+cJ$, and both
  coefficients vanish.~\qed\end{proof}
 \step{s414srgcn}{$cn=(k+2)(k-\eta)$.}
 \begin{proof}\pf\ Apply \stepref{s414srgtheta} to $\Zj$: $Q\Zj=(k+2)\Zj$ and
  $(A-\eta I)\Zj=(k-\eta)\Zj$ by \stepref{s414srgmat}, so
  $(k-\eta)(k+2)\Zj=cJ\Zj=cn\Zj$.~\qed\end{proof}
 \pflet{$E=Q-\frac{k+2}nJ$.}
 \step{s414srgPJ}{$PJ=JP=0$, and $E=Q(I-\frac1nJ)$.}
 \begin{proof}\pf\ $QJ=JQ=(k+2)J$ and $J^2=nJ$ by \stepref{s414srgmat}.~\qed\end{proof}
 \step{s414srgidem}{$E^2=(\eta+2)E$.}
 \begin{proof}\pf\
  \[
  \begin{split}
   E^2&=PQ-\tfrac{k+2}nPJ=PQ=Q^2-\tfrac{(k+2)^2}nJ\\
      &=(A-\eta I)Q+(\eta+2)Q-\tfrac{(k+2)^2}nJ
       =(\eta+2)Q+\Bigl(c-\tfrac{(k+2)^2}n\Bigr)J\\
      &=(\eta+2)Q-(\eta+2)\tfrac{k+2}nJ=(\eta+2)E,
  \end{split}
  \]
  where the first line uses \stepref{s414srgPJ} and $JQ=(k+2)J$, the second
  \stepref{s414srgtheta}, and the third \stepref{s414srgcn}, which gives
  $cn-(k+2)^2=-(k+2)(\eta+2)$.~\qed\end{proof}
 \step{s414srgrkP}{$(\eta+2)\rk E=2n-k-2$.}
 \begin{proof}\pf\ $\eta+2\ge2$ by \stepref{s414srgtheta}, so $E/(\eta+2)$ is idempotent by
  \stepref{s414srgidem}, and its rank is its trace (\path{LinearMap.IsProj.trace} and
  \path{LinearMap.isProj_range_iff_isIdempotentElem}). Further $\tr E=\tr Q-\frac{k+2}n\tr J=2n-(k+2)$,
  since $\tr A=0$.~\qed\end{proof}
 \step{s414srgrkQ}{$\rk Q=\rk E+1$.}
 \begin{proof}\pf\ The image of $E=Q(I-\frac1nJ)$ lies in that of $Q$, and so does
  $\Zj=Q\Zj/(k+2)$; conversely $Q=E+\frac{k+2}nJ$ has image in $\opk{im}E+\re\Zj$. Finally
  $\Zj\notin\opk{im}E$: $JP=0$ (\stepref{s414srgPJ}) gives $(\Zj,E\Zy)=0$ for every $\Zy$, while
  $(\Zj,\Zj)=n>0$.~\qed\end{proof}
 \qedstep
 \begin{proof}\pf\ Assertion 1 is \stepref{s414srgnc}, \stepref{s414srgkn}, \stepref{s414srgak},
  \stepref{s414srgck} and \stepref{s414srgc}, assertion 2 is \stepref{s414srgunion} and assertion 3 is
  \stepref{s414srglin}. For assertion 4, $\eta+2=a-c+4$, and \stepref{s414srgrkP} and
  \stepref{s414srgrkQ} give $(\eta+2)(\rk Q-1)=2n-k-2$.~\qed\end{proof}
\end{proof}

\begin{namedresult}{Computation}{C18}\label{comp:C18}
 \lean{Computations.C18.srg_parameters}
 \uses{}
 \assume{$n,k,a,c,\rho$ are integers with $1\le k\le n-2$, $n\le36$, $k\le16$, $0\le a\le k-1$,
 $1\le c\le k$ and $0\le\rho\le8$, such that
 \begin{pfenum}
  \item $k(k-a-1)=c(n-k-1)$,
  \item $k-2a+c=4$,
  \item $2k-c\le n-2$, and
  \item $(a-c+4)(\rho-1)=2n-k-2$.
 \end{pfenum}}
 \prove{$(n,k,a,c,\rho)$ is one of the following 15 tuples:
 $(2m,2m-2,2m-4,2m-2,m+1)$ for $2\le m\le7$; $(9,4,1,2,5)$; $(10,6,3,4,5)$;
 $(10,3,0,1,6)$; $(15,8,4,4,6)$; $(16,10,6,6,6)$; $(16,6,2,2,7)$; $(21,10,5,4,7)$;
 $(27,16,10,8,7)$; and $(28,12,6,4,8)$.}
 This is Lemma P of \path{cgss/prerequisites.md} (SR5), with $\rho=1+m_1$. The first six are
 the parameters of $CP(m)$, the next two those of $L_2(3)$ and $T(5)$.
\end{namedresult}
\begin{proof}
 \pfsketch\ This is a finite search. The Lean proof is one \texttt{Bool} search over the
 valencies with a soundness lemma, as for the checker \path{Computations.allSRGParameters} of
 \path{LineGraphs/Computations/BoundedSearch.lean}, whose pruning ($a<k$, $0<c\le k$, $k+1\le n$)
 matches our bounds; \texttt{interval\_cases} would make about $1.5\cdot10^5$ goals.
 \qedstep
 \begin{proof}\pf\ For each $k\le16$, $a<k$ and $1\le c\le k$, condition 1 determines $n$, and
  condition 4 determines $\rho$, since $a-c+4=0$ would give $2n=k+2$, against $k+2\le n$ and
  $n\ge3$; the search keeps the tuples satisfying all four
  conditions and the bounds. It finds the 15 tuples, and no other (SageMath,
  \path{cgss/certificates/c18_parameter_search.sage}).
  Without condition 3 the tuple $(21,16,12,12,7)$ would also survive.~\qed\end{proof}
\end{proof}

\begin{namedresult}{Lemma}{(strongly regular line graphs)}\label{lem:414-linegraph}
 \lean{SimpleGraph.IsSRGWith.nonempty_iso_lineGraph_complete_or_completeBipartite}
 \uses{lem:4-LG1,lem:414-iso,facts:414}
 \assume{\begin{pfenum}
  \item $X$ is a connected graph, strongly regular with parameters $(n,k,a,c)$.
  \item $\det(2I+A(X))=0$.
  \item $Y$ is a graph and $X\cong L(Y)$.
 \end{pfenum}}
 \prove{$X\cong L(K_m)$ or $X\cong L(K_{m,m})$ for some $m\ge2$.}
 \new\ This is the line graph case of Theorem 4.13, stated as a result so that we can cite it;
 its argument is that of the proof of Theorem 4.13, through Lemma LG1.
\end{namedresult}
\begin{proof}
 \pfsketch\ We will remove the isolated vertices of $Y$ and apply Lemma LG1; the eigenvalue $-2$
 excludes complete graphs and $C_5$.
 \pflet{$Y'$ be $Y$ with its isolated vertices removed.}
 \step{s414lgsame}{$L(Y')=L(Y)$.}
 \begin{proof}\pf\ Removing isolated vertices keeps the edges and their
  intersections.~\qed\end{proof}
 \step{s414lgedge}{$Y'$ has an edge.}
 \begin{proof}\pf\ $X$ has a vertex, being connected, so $L(Y)\cong X$ (\asm[3]{0}) has one; it is
  an edge of $Y$, and of $Y'$ by \stepref{s414lgsame}.~\qed\end{proof}
 \step{s414lgconn}{$Y'$ is connected.}
 \begin{proof}\pf\ Every vertex of $Y'$ lies on an edge, and a path in the connected graph
  $L(Y')\cong X$ (\stepref{s414lgsame}, \asm[3]{0}) between two edges gives a walk in $Y'$ between
  their ends.~\qed\end{proof}
 \step{s414lgcases}{$Y'$ is a cycle, or $Y'\cong K_m$, $Y'\cong K_{q,q}$ or $Y'\cong K_{1,q}$
  for some $m$ or $q$.}
 \begin{proof}\pf\ Lemma~\ref{lem:4-LG1}, with \stepref{s414lgedge} and \stepref{s414lgconn},
  since $L(Y')$ is strongly regular by \asm[1]{0}, \asm[3]{0}, \stepref{s414lgsame} and
  Lemma~\ref{lem:414-iso}, assertion 2.~\qed\end{proof}
 \step{s414lgKn}{$X$ is not complete.}
 \begin{proof}\pf\ If $X$ is complete on $n\ge1$ vertices, then $2I+A=I+J$, which is positive
  definite since $(\Zz,(I+J)\Zz)=|\Zz|^2+(\Zz,\Zj)^2$; so its determinant is not 0, against
  \asm[2]{0}.~\qed\end{proof}
 \step{s414lgCfivedet}{$\det(2I+A(C_5))=4$.}
 \begin{proof}\pf\ \texttt{decide} on the $5\times5$ integer matrix.~\qed\end{proof}
 \step{s414lgC5}{$X\not\cong C_5$.}
 \begin{proof}\pf\ An isomorphism $X\to C_5$ conjugates $2I+A(X)$ into $2I+A(C_5)$ by a
  permutation matrix, which keeps the determinant; apply \stepref{s414lgCfivedet} and
  \asm[2]{0}.~\qed\end{proof}
 \step{s414lgCp}{$X\not\cong C_p$ for $p\ge6$.}
 \begin{proof}\pf\ In $C_p$ with $p\ge6$ there are non-adjacent vertices at distance 2 and at
  distance 3, with 1 and 0 common neighbours; so $C_p$ is not strongly regular, against
  \asm[1]{0}.~\qed\end{proof}
 \step{s414lgcycle}{\case{$Y'$ is a cycle $C_p$.}\prove{$X\cong L(K_{2,2})$.}}
 \begin{proof}\pf\ $L(C_p)\cong C_p$, and $p\ge3$. The cases $p=3$ (complete), $p=5$ and $p\ge6$
  are excluded by \stepref{s414lgKn}, \stepref{s414lgC5} and \stepref{s414lgCp}. So $p=4$, and
  $C_4\cong L(C_4)=L(K_{2,2})$.~\qed\end{proof}
 \step{s414lgstar}{\case{$Y'\cong K_{1,q}$.}\prove{This case does not occur.}}
 \begin{proof}\pf\ Any two edges of $K_{1,q}$ meet, so $L(K_{1,q})\cong K_q$ is complete, against
  \stepref{s414lgKn}.~\qed\end{proof}
 \qedstep
 \begin{proof}\pf\ By \stepref{s414lgcases} the cases \stepref{s414lgcycle},
  \stepref{s414lgstar}, $Y'\cong K_m$ and $Y'\cong K_{q,q}$ are exhaustive. In the last two,
  $X\cong L(K_m)$ or $L(K_{q,q})$ by \asm[3]{0} and \stepref{s414lgsame}. Further $m\ge2$, since
  $L(K_m)$ has no vertex for $m\le1$ while the connected graph $X$ has one; and $q\ge2$, since
  $L(K_{0,0})$ has no vertex and $L(K_{1,1})$ is complete (\stepref{s414lgKn}).~\qed\end{proof}
\end{proof}

\subsubsection*{The root system of a representation}

\begin{namedresult}{Lemma}{(the lattice roots)}\label{lem:414-lattice}
 \lean{SimpleGraph.IsRepresentation.latticeRoots}
 \lean{SimpleGraph.IsRepresentation.isStarClosed_latticeRoots}
 \lean{SimpleGraph.IsRepresentation.isIndecomposable_latticeRoots}
 \lean{SimpleGraph.IsRepresentation.span_latticeRoots}
 \lean{SimpleGraph.IsRepresentation.isRepresentedIn_latticeRoots}
 \uses{def:represented,def:line-system,def:indecomposable,def:2.1,lem:4-lines,lem:sec3-connected-span,facts:414}
 \assume{$X$ is a connected graph with a representation $(\Zx_u)$ in a finite-dimensional real
 inner product space $V$; $\Lambda$ is the $\ints$-span and $W$ the span of the $\Zx_u$; and
 $\Phi_R=\{\Zy\in \Lambda:(\Zy,\Zy)=2\}$.}
 \prove{\begin{pfenum}
  \item $(\Zy,\Zz)\in\ints$ for all $\Zy,\Zz\in \Lambda$.
  \item $\Phi_R$ is the set of vectors of a line system $R$ in $V$, which is nonempty,
   star-closed and indecomposable and spans $W$.
  \item $X$ is represented in $R$ by $(\Zx_u)$.
 \end{pfenum}}
 \new\ \path{cgss/prerequisites.md} (LN3) uses the star closure of the $\Zx_u$. The larger set
 $\Phi_R$ is star-closed as well, and it lies in $\Lambda$ by definition, which is what the apex
 argument needs; we do not need $R\sbs\typ E_8$.
\end{namedresult}
\begin{proof}
 \pfsketch\ Inner products in $\Lambda$ are integers, so norms and Cauchy--Schwarz do the rest.
 \step{s414latint}{Assertion 1 holds.}
 \begin{proof}\pf\ $\Zy=\sum_u y_u\Zx_u$ and $\Zz=\sum_w z_w\Zx_w$ with integers $y_u,z_w$, and
  $(\Zy,\Zz)=\sum_{u,w}y_uz_w(\Zx_u,\Zx_w)$, where each $(\Zx_u,\Zx_w)\in\{0,1,2\}$ by
  Definition~\ref{def:represented}.~\qed\end{proof}
 \step{s414latvals}{$\Phi_R=-\Phi_R$, and $(\Zy,\Zz)\in\{-2,-1,0,1,2\}$ for
  $\Zy,\Zz\in\Phi_R$.}
 \begin{proof}\pf\ $\Lambda=-\Lambda$. By \path{abs_real_inner_le_norm},
  $|(\Zy,\Zz)|\le|\Zy||\Zz|=2$, and $(\Zy,\Zz)$ is an integer by
  \stepref{s414latint}.~\qed\end{proof}
 \step{s414latmem}{$\Zx_u\in\Phi_R$ for every $u$, and $\Phi_R\sbs \Lambda\sbs W$.}
 \begin{proof}\pf\ $(\Zx_u,\Zx_u)=2$ by Definition~\ref{def:represented}; $\Lambda\sbs W$ since integer
  combinations are combinations.~\qed\end{proof}
 \step{s414latbasis}{\pick\ vertices $u_1,\dots,u_\rho$ such that $\Zx_{u_1},\dots,\Zx_{u_\rho}$
  form a basis of $W$.}
 \begin{proof}\pf\ \path{exists_linearIndependent}, since the $\Zx_u$ span $W$ and $W$ is
  finite-dimensional.~\qed\end{proof}
 \step{s414latinj}{The map $\iota:\Zy\mapsto((\Zy,\Zx_{u_i}))_{i\le\rho}$ is injective on $W$.}
 \begin{proof}\pf\ It is linear, and a vector of $W$ orthogonal to the basis of
  \stepref{s414latbasis} is orthogonal to $W$, so to itself, and is 0.~\qed\end{proof}
 \step{s414latfin}{$\Phi_R$ is finite.}
 \begin{proof}\pf\ By \stepref{s414latvals} and \stepref{s414latmem}, $\iota$ maps $\Phi_R\sbs W$
  into the finite set $\{-2,\dots,2\}^\rho$; apply \stepref{s414latinj}.~\qed\end{proof}
 \step{s414latstar}{$\Phi_R$ is star-closed.}
 \begin{proof}\pf\ For $\Zy,\Zz\in\Phi_R$ with $(\Zy,\Zz)=-1$, $\Zy+\Zz\in \Lambda$ and
  $(\Zy+\Zz,\Zy+\Zz)=2-2+2=2$; this is Definition~\ref{def:2.1}.~\qed\end{proof}
 \step{s414latspan}{$\Phi_R$ spans $W$.}
 \begin{proof}\pf\ By \stepref{s414latmem}, $\{\Zx_u\}\sbs\Phi_R\sbs W$, and the $\Zx_u$ span
  $W$.~\qed\end{proof}
 \step{s414latind}{$R$ is nonempty and indecomposable.}
 \begin{proof}\pf\ Take $B=\{\Zx_u\}$ in Lemma~\ref{lem:sec3-connected-span}: $B$ is nonempty
  since $X$ has a vertex, it spans the span of $\Phi_R$ by \stepref{s414latspan}, and distinct
  members of $B$ have nonzero inner product exactly when their vertices are adjacent
  (Lemma~\ref{lem:4-lines}, assertion 1, and Definition~\ref{def:represented}), so its graph is a
  copy of the connected graph $X$.~\qed\end{proof}
 \qedstep
 \begin{proof}\pf\ Assertion 1 is \stepref{s414latint}. By \stepref{s414latvals},
  \stepref{s414latfin} and Definition~\ref{def:line-system}, $\Phi_R$ is the set of vectors of a
  line system $R$; \stepref{s414latstar}, \stepref{s414latspan} and \stepref{s414latind} give the
  rest of assertion 2. Assertion 3 is \stepref{s414latmem} with
  Definition~\ref{def:represented}.~\qed\end{proof}
\end{proof}

\begin{namedresult}{Lemma}{(the apex)}\label{lem:414-apex}
 \lean{SimpleGraph.IsRepresentation.apex}\lean{SimpleGraph.IsRepresentation.inner_apex}
 \lean{SimpleGraph.IsRepresentation.inner_apex_self}
 \lean{SimpleGraph.IsRepresentation.inner_apex_mem_int}
 \lean{SimpleGraph.IsRepresentation.apex_mem_span}
 \uses{def:represented}
 \assume{$X$ is a $k$-regular graph on $n\ge1$ vertices with a representation $(\Zx_u)$ in a
 real inner product space $V$; $\Lambda$ is the $\ints$-span and $W$ the span of the $\Zx_u$; and
 $\Zw=\frac1{k+2}\sum_u\Zx_u$.}
 \prove{\begin{pfenum}
  \item $(\Zw,\Zx_v)=1$ for every vertex $v$.
  \item $(\Zw,\Zw)=n/(k+2)$.
  \item $(\Zw,\Zy)\in\ints$ for every $\Zy\in \Lambda$.
  \item $\Zw\in W$.
 \end{pfenum}}
 \new\ Step $\langle1\rangle5$ of \path{cgss/prerequisites.md}, Section 7.3, asks for integral
 inner products with the roots of $R$ only; we prove them on the whole $\ints$-span, and add
 assertion 4, which the switching from the apex needs.
\end{namedresult}
\begin{proof}
 \pfsketch\ This is step $\langle1\rangle5$ of \path{cgss/prerequisites.md}, Section 7.3; the
 row sums of $2I+A$ are $k+2$.
 \step{s414apexone}{Assertion 1 holds.}
 \begin{proof}\pf\ $(\Zw,\Zx_v)=\frac1{k+2}\sum_u(\Zx_u,\Zx_v)=\frac1{k+2}(2+k)$, since by
  Definition~\ref{def:represented} the terms are 2 for $u=v$, 1 for the $k$ neighbours of $v$ and 0
  otherwise.~\qed\end{proof}
 \step{s414apextwo}{Assertion 2 holds.}
 \begin{proof}\pf\ $(\Zw,\Zw)=\frac1{k+2}\sum_v(\Zw,\Zx_v)=\frac n{k+2}$ by
  \stepref{s414apexone}.~\qed\end{proof}
 \step{s414apexthree}{Assertion 3 holds.}
 \begin{proof}\pf\ For $\Zy=\sum_vz_v\Zx_v$ with integers $z_v$, $(\Zw,\Zy)=\sum_vz_v$ by
  \stepref{s414apexone}.~\qed\end{proof}
 \qedstep
 \begin{proof}\pf\ \stepref{s414apexone}, \stepref{s414apextwo}, \stepref{s414apexthree}, and
  $\Zw$ is a combination of the $\Zx_u$.~\qed\end{proof}
\end{proof}

\begin{namedresult}{Lemma}{(integrality)}\label{lem:414-dual}
 \lean{InnerProductSpace.det_gram_mul_inner_self_mem_int}
 \uses{facts:414}
 \assume{$W$ is a real inner product space with a basis $\Zb_1,\dots,\Zb_t$ whose Gram matrix
 $B$ has integer entries, and $\Zw\in W$ has $(\Zw,\Zb_i)\in\ints$ for every $i$.}
 \prove{$\det(B)\,(\Zw,\Zw)$ is an integer.}
 \new\ Step $\langle1\rangle6$ of \path{cgss/prerequisites.md}, Section 7.3, writes
 $(\Zw,\Zw)=\Za^TB^{-1}\Za$ and cites LA6, which lists tools only; we state the integrality as a
 result and prove it with the adjugate.
\end{namedresult}
\begin{proof}
 \pfsketch\ We will write $(\Zw,\Zw)$ with the inverse of $B$, and clear the denominator with
 the adjugate.
 \pflet{$\Za\in\ints^t$ with $\Za(i)=(\Zw,\Zb_i)$, and $\beta\in\re^t$ with
 $\Zw=\sum_i\beta(i)\Zb_i$.}
 \step{s414dualeq}{$\Za=B\beta$ and $(\Zw,\Zw)=\beta^T\Za$.}
 \begin{proof}\pf\ $\Za(i)=\sum_j\beta(j)(\Zb_j,\Zb_i)$ and
  $(\Zw,\Zw)=\sum_i\beta(i)(\Zw,\Zb_i)$, by bilinearity.~\qed\end{proof}
 \step{s414dualinv}{$\det B\ne0$.}
 \begin{proof}\raggedright\pf\ The $\Zb_i$ are linearly independent; apply
  \path{Matrix.det_gram_ne_zero_iff_linearIndependent}.~\qed\end{proof}
 \step{s414dualadj}{$\det(B)(\Zw,\Zw)=\Za^T\opk{adj}(B)\,\Za$.}
 \begin{proof}\pf\
  $\det(B)(\Zw,\Zw)=\det(B)\,\Za^T\beta=\det(B)\,\Za^TB^{-1}\Za=\Za^T\opk{adj}(B)\,\Za$, by
  \stepref{s414dualeq}, then $\beta=B^{-1}\Za$ with \stepref{s414dualinv}, then
  $B^{-1}=\det(B)^{-1}\opk{adj}(B)$ (\path{Matrix.inv_def}).~\qed\end{proof}
 \qedstep
 \begin{proof}\pf\ The entries of $\opk{adj}(B)$ are signed minors of the integer matrix $B$, and
  $\Za$ is integral; apply \stepref{s414dualadj}.~\qed\end{proof}
\end{proof}

\begin{namedresult}{Computation}{C21}\label{comp:C21}
 \lean{LineSystem.E₆.simpleRoot}\lean{LineSystem.E₆.simpleRoot_mem}
 \lean{LineSystem.E₆.gram_simpleRoot}\lean{LineSystem.E₆.det_gram_simpleRoot}
 \lean{LineSystem.E₆.simpleBasis}\lean{LineSystem.E₆.gram_simpleRoot_eq_cartan}
 \lean{LineSystem.E₇.simpleRoot}\lean{LineSystem.E₇.simpleRoot_mem}
 \lean{LineSystem.E₇.gram_simpleRoot}\lean{LineSystem.E₇.det_gram_simpleRoot}
 \lean{LineSystem.E₇.simpleBasis}\lean{LineSystem.E₇.gram_simpleRoot_eq_cartan}
 \lean{LineSystem.E₈.simpleRoot}\lean{LineSystem.E₈.simpleRoot_mem}
 \lean{LineSystem.E₈.gram_simpleRoot}\lean{LineSystem.E₈.det_gram_simpleRoot}
 \lean{LineSystem.E₈.simpleBasis}\lean{LineSystem.E₈.gram_simpleRoot_eq_cartan}
 \lean{Computations.C21.e6_gram_det}\lean{Computations.C21.e7_gram_det}
 \lean{Computations.C21.e8_gram_det}
 \uses{def:3.2,def:3.3,lem:sec3-spans}
 \assume{$t\in\{6,7,8\}$.}
 \prove{There are roots $\Zs_1,\dots,\Zs_t\in\Phi(\typ E_t)$ which form a basis of the span of
 $\Phi(\typ E_t)$, and whose Gram matrix $B_t$ has integer entries and determinant 3, 2 and 1
 for $t=6,7,8$. Up to a relabelling, $B_t$ is Mathlib's \texttt{CartanMatrix.E}~$t$.}
 This is LN7 of \path{cgss/prerequisites.md}.
\end{namedresult}
\begin{proof}
 \pfsketch\ The roots are those of the SageMath certificate. Lean checks their membership, their
 Gram matrix and a scaled $LDL^T$ factorization of it by \texttt{decide +kernel}; see
 \path{LineGraphs/Computations/C21E6.lean} and its analogues.
 \qedstep
 \begin{proof}\pf\ \leancite{LineSystem.E₆.simpleRoot_mem}{\texttt{LineSystem.E\ensuremath{{}_6}.simpleRoot\_mem}},
  \leancite{LineSystem.E₆.gram_simpleRoot}{\texttt{E\ensuremath{{}_6}.gram\_simpleRoot}} and
  \leancite{LineSystem.E₆.det_gram_simpleRoot}{\texttt{E\ensuremath{{}_6}.det\_gram\_simpleRoot}}
  give the roots, their integer Gram matrix and its determinant; the Gram determinant is not 0,
  so the $t$ roots are independent, and they span since the span has dimension $t$
  (Lemma~\ref{lem:sec3-spans}), which is
  \leancite{LineSystem.E₆.simpleBasis}{\texttt{E\ensuremath{{}_6}.simpleBasis}}. The same
  declarations for $\typ E_7$, $\typ E_8$ give the other two cases.~\qed\end{proof}
\end{proof}

\begin{namedresult}{Computation}{C20 at the apex}\label{comp:414-Nc}
 \lean{LineSystem.E₆.c_mem}\lean{LineSystem.E₇.c_mem}\lean{LineSystem.E₈.c_mem}
 \lean{LineSystem.E₆.N_c}\lean{LineSystem.E₇.N_c}\lean{LineSystem.E₈.N_c}
 \lean{LineSystem.E₆.card_CEdge}\lean{LineSystem.E₇.card_CEdge}\lean{LineSystem.E₈.card_CEdge}
 \uses{def:3.2,def:3.3,comp:C20}
 \assume{$t\in\{6,7,8\}$, and $N_\Zc=\{\Zx\in\Phi(\typ E_t):(\Zx,\Zc)=1\}$.}
 \prove{$\Zc\in\Phi(\typ E_t)$, $N_\Zc=\{\Ze_P:P\in\Ecal_t\}\cup\{\Zc-\Ze_P:P\in\Ecal_t\}$, and
 $|\Ecal_t|$ is 10, 16 and 28 for $t=6,7,8$.}
 This is LN6 of \path{cgss/prerequisites.md}, at the root $\Zc$.
\end{namedresult}
\begin{proof}
 \pfsketch\ Lean identifies the pairs of $N_\Zc$ that Computation C20 lists with the members of
 $\Ecal_t$, in \path{LineGraphs/Computations/C20Apex.lean}.
 \qedstep
 \begin{proof}\pf\ \leancite{LineSystem.E₆.c_mem}{\texttt{LineSystem.E\ensuremath{{}_6}.c\_mem}},
  \leancite{LineSystem.E₆.N_c}{\texttt{E\ensuremath{{}_6}.N\_c}} and
  \leancite{LineSystem.E₆.card_CEdge}{\texttt{E\ensuremath{{}_6}.card\_CEdge}}, and their analogues for
  $\typ E_7$ and $\typ E_8$; in Lean the coordinates are numbered from 0, and $\Ecal_t$ is the
  subtype \leancite{LineSystem.E₆.CEdge}{\texttt{E\ensuremath{{}_6}.CEdge}} of ordered pairs
  $i<j$.~\qed\end{proof}
\end{proof}

\begin{namedresult}{Lemma}{(transitivity on roots)}\label{lem:414-transitive}
 \lean{LineSystem.E₆.exists_isometry_apply_eq_c}\lean{LineSystem.E₇.exists_isometry_apply_eq_c}
 \lean{LineSystem.E₈.exists_isometry_apply_eq_c}
 \uses{facts:414,comp:414-Nc,comp:C4,lem:3.4,def:3.3,def:indecomposable,facts:inner}
 \assume{$t\in\{6,7,8\}$ and $\Zr\in\Phi(\typ E_t)$.}
 \prove{There is a linear isometry $\psi$ of $\re^8$ with $\psi(\Phi(\typ E_t))=\Phi(\typ E_t)$ and
 $\psi(\Zr)=\Zc$.}
 \new\ This replaces LN5 of \path{cgss/prerequisites.md}, a transitivity by reflections, by the
 transitivity on ordered stars that the Lean of Computation C4 proves.
\end{namedresult}
\begin{proof}
 \pfsketch\ Both $\Zr$ and $\Zc$ lie in ordered stars, and the stars form one orbit.
 \step{s414trr}{There is $\Zs\in\Phi(\typ E_t)$ with $(\Zr,\Zs)=-1$.}
 \begin{proof}\pf\ $\typ E_t$ is indecomposable (Computation~\ref{comp:C4}) and has more than one
  line (Lemma~\ref{lem:3.4}). If every root other than $\pm\Zr$ were orthogonal to $\Zr$, then
  $\{\pm\Zr\}$ and the other roots would split $\Phi(\typ E_t)$ orthogonally into nonempty parts
  (the definition of indecomposability). So $(\Zr,\Zs')\ne0$ for some $\Zs'\ne\pm\Zr$, and then
  $(\Zr,\Zs')=\pm1$ by \eqref{eq:endpoints}; take $\Zs=\mp\Zs'$.~\qed\end{proof}
 \step{s414trc}{$\Zc$ and $-\Ze_1-\Ze_2$ lie in $\Phi(\typ E_t)$, and
  $(\Zc,-\Ze_1-\Ze_2)=-1$.}
 \begin{proof}\pf\ $\Zc\in\Phi(\typ E_t)$ by Computation~\ref{comp:414-Nc}; $-\Ze_1-\Ze_2$ lies in
  $\Phi(\typ D_8)$ and has $x_6=x_7=x_8=0$, so it lies in $\Phi(\typ E_6)\sbs\Phi(\typ E_7)$ by
  Definition~\ref{def:3.3}; and $(\Zc,\Ze_1+\Ze_2)=\frac12+\frac12$.~\qed\end{proof}
 \qedstep
 \begin{proof}\pf\ The transitivity on ordered stars of Computation~\ref{comp:C4} (in
  Lean; see the facts at the start of this section), applied to the star
  $(\Zr,\Zs)$ of \stepref{s414trr} and to the star $(\Zc,-\Ze_1-\Ze_2)$ of \stepref{s414trc},
  gives isometries $g_1,g_2$ preserving $\Phi(\typ E_t)$ with $g_1(\Za_0)=\Zr$ and
  $g_2(\Za_0)=\Zc$; take $\psi=g_2g_1^{-1}$.~\qed\end{proof}
\end{proof}

\begin{namedresult}{Lemma}{(switching from the apex)}\label{lem:414-transversal}
 \lean{SimpleGraph.IsRepresentation.nonempty_iso_lineGraph_switch}
 \uses{def:represented,def:414-switch,def:2.2,thm:3.6,lem:4-lines,lem:4-transport,lem:414-lattice,lem:414-apex,lem:414-transitive,comp:414-Nc,facts:414}
 \assume{\begin{pfenum}
  \item $X$ is a connected $k$-regular graph on $n$ vertices with a representation $(\Zx_u)$ in a
   finite-dimensional real inner product space $V$.
  \item $\Lambda$, $W$ and $R$ are as in Lemma~\ref{lem:414-lattice}, and $\Zw$ is as in
   Lemma~\ref{lem:414-apex}.
  \item $t\in\{6,7,8\}$, and $f:\Phi(R)\to\Phi(\typ E_t)$ is a bijection that preserves inner
   products.
  \item $(\Zw,\Zw)=2$ and $n=|\Ecal_t|$.
 \end{pfenum}}
 \prove{There are a bijection $\pi:V(X)\to\Ecal_t$ and a set $H\sbs\Ecal_t$ such that $\pi$ is
 an isomorphism from $X$ to $L(\Delta_t)^H$.}
 \new\ These are steps $\langle1\rangle7$ to $\langle1\rangle9$ of
 \path{cgss/prerequisites.md}, Section 7.3; we read the switched graph off the inner products
 directly, without SW4.
\end{namedresult}
\begin{proof}
 \pfsketch\ The apex is a root, the $\Zx_u$ are a transversal of the pairs of $N_\Zw$, and an
 isometry takes $\Zw$ to $\Zc$, where the pairs are indexed by $\Ecal_t$.
 \step{s414trmax}{$R$ is maximal in $W$.}
 \begin{proof}\pf\ $R$ is nonempty and indecomposable and spans $W$ (Lemma~\ref{lem:414-lattice});
  it is isomorphic to $\typ E_t$ by \asm[3]{0}, which is in the list of
  Theorem~\ref{thm:3.6}.~\qed\end{proof}
 \step{s414trw}{$\Zw\in\Phi(R)$.}
 \begin{proof}\pf\ Lemma~\ref{lem:414-apex} applies, since $n=|\Ecal_t|\ge1$ by \asm[4]{0} and
  Computation~\ref{comp:414-Nc}; so $\Zw\in W$, and $(\Zw,\Zw)=2$ by \asm[4]{0}. For
  $\Zy\in\Phi(R)\sbs \Lambda$, $(\Zw,\Zy)$ is an integer (Lemma~\ref{lem:414-apex}) with
  $|(\Zw,\Zy)|\le|\Zw||\Zy|=2$ (\path{abs_real_inner_le_norm}); apply \stepref{s414trmax} and
  Definition~\ref{def:2.2}.~\qed\end{proof}
 \step{s414trg}{There is a bijection $g:\Phi(R)\to\Phi(\typ E_t)$ that preserves inner products,
  with $g(\Zw)=\Zc$.}
 \begin{proof}\pf\ $f(\Zw)\in\Phi(\typ E_t)$ by \stepref{s414trw}; Lemma~\ref{lem:414-transitive}
  gives $\psi$ with $\psi(f(\Zw))=\Zc$, and $g=\psi\circ f$ will do.~\qed\end{proof}
 \pflet{$\Zy_u=g(\Zx_u)$.}
 \step{s414trrep}{$(\Zy_u)$ represents $X$ in $\typ E_t$, and $(\Zy_u,\Zc)=1$ for every $u$.}
 \begin{proof}\pf\ $X$ is represented in $R$ by $(\Zx_u)$ (Lemma~\ref{lem:414-lattice}); apply
  Lemma~\ref{lem:4-transport} to $g$. Further
  $(\Zy_u,\Zc)=(g(\Zx_u),g(\Zw))=(\Zx_u,\Zw)=1$ by Lemma~\ref{lem:414-apex}.~\qed\end{proof}
 \step{s414trpick}{\pick\ $\pi(u)\in\Ecal_t$ for each $u$, and $D\sbs V(X)$, such that
  $\Zy_u=\Zc-\Ze_{\pi(u)}$ for $u\in D$ and $\Zy_u=\Ze_{\pi(u)}$ for $u\notin D$.}
 \begin{proof}\pf\ $\Zy_u\in N_\Zc$ by \stepref{s414trrep}; apply
  Computation~\ref{comp:414-Nc}.~\qed\end{proof}
 \step{s414trip}{For $P,P'\in\Ecal_t$, $(\Ze_P,\Ze_{P'})=(\Zc-\Ze_P,\Zc-\Ze_{P'})=|P\cap P'|$
  and $(\Ze_P,\Zc-\Ze_{P'})=1-|P\cap P'|$.}
 \begin{proof}\pf\ Expand, with $(\Zc,\Zc)=2$, $(\Zc,\Ze_i)=\frac12$ and
  $(\Ze_i,\Ze_j)=[i=j]$.~\qed\end{proof}
 \step{s414trinj}{$\pi$ is injective.}
 \begin{proof}\pf\ If $u\ne v$ and $\pi(u)=\pi(v)=P$, then $\Zy_u\ne\Zy_v$
  (Lemma~\ref{lem:4-lines}, assertion 1), so one is $\Ze_P$ and the other $\Zc-\Ze_P$, and their
  inner product is $1-2=-1$ by \stepref{s414trip}; this is against
  Definition~\ref{def:represented} and \stepref{s414trrep}.~\qed\end{proof}
 \step{s414trbij}{$\pi$ is a bijection.}
 \begin{proof}\pf\ \stepref{s414trinj}, and $|V(X)|=n=|\Ecal_t|$ by \asm[4]{0}.~\qed\end{proof}
 \pflet{$H=\pi(D)$.}
 \step{s414tradj}{Distinct $u,v$ are adjacent in $X$ if and only if $\pi(u)$ and $\pi(v)$ are
  adjacent in $L(\Delta_t)^H$.}
 \begin{proof}\pf\ By \stepref{s414trinj}, $|\pi(u)\cap\pi(v)|\le1$, and $u\in D$ exactly when
  $\pi(u)\in H$. By \stepref{s414trrep} and \stepref{s414trip}, $(\Zy_u,\Zy_v)=1$ exactly when
  $\pi(u),\pi(v)$ meet and both or neither lie in $H$, or are disjoint and exactly one lies in
  $H$; this is adjacency in $L(\Delta_t)^H$ (Definition~\ref{def:414-switch}).~\qed\end{proof}
 \qedstep
 \begin{proof}\pf\ \stepref{s414trbij} and \stepref{s414tradj}.~\qed\end{proof}
\end{proof}

\subsubsection*{Regular switchings}

\begin{namedresult}{Lemma}{(degrees in a switched line graph)}\label{lem:414-degree}
 \lean{SimpleGraph.card_adj_lineGraph_switch_of_complete}
 \uses{def:414-switch}
 \assume{\begin{pfenum}
  \item $Y$ is a graph, $M\sbs V(Y)$ has $m\ge2$ elements, and every 2-subset of $M$ is an edge
   of $Y$; $K$ is the set of these edges.
  \item $H\sbs E(Y)$, $h=|H\cap K|$, and $\deg_{H\cap K}$ is taken on $M$.
  \item $e=\{i,j\}\in K$ and $s=\deg_{H\cap K}(i)+\deg_{H\cap K}(j)$.
 \end{pfenum}}
 \prove{The number of neighbours of $e$ in $L(Y)^H$ that lie in $K$ is
 $2s-h-3+\binom{m-2}2$ if $e\in H$, and $2(m-2)+h-2s$ if $e\notin H$.}
 \new\ This replaces SW3 of \path{cgss/prerequisites.md} in step $\langle1\rangle10$ of its
 Section 7.3 by a direct count.
\end{namedresult}
\begin{proof}
 \pfsketch\ We will count the members of $K$ that meet $e$ and those that miss it, inside and
 outside $H$.
 \step{s414degmeet}{The members of $H\cap K$ other than $e$ that meet $e$ number $s-2$ if $e\in H$
  and $s$ if $e\notin H$.}
 \begin{proof}\pf\ No member of $K$ other than $e$ contains both $i$ and $j$, so the members of
  $H\cap K$ meeting $e$ number $\deg_{H\cap K}(i)+\deg_{H\cap K}(j)$, counting $e$ twice when
  $e\in H$.~\qed\end{proof}
 \step{s414degall}{The members of $K$ other than $e$ that meet $e$ number $2(m-2)$, and those
  disjoint from $e$ number $\binom{m-2}2$.}
 \begin{proof}\pf\ They are $\{i,l\}$ and $\{j,l\}$ with $l\in M\setminus e$, and the 2-subsets
  of $M\setminus e$.~\qed\end{proof}
 \step{s414degin}{\case{$e\in H$.}\prove{The count is $2s-h-3+\binom{m-2}2$.}}
 \begin{proof}\pf\ By Definition~\ref{def:414-switch}, the neighbours of $e$ in $K$ are the
  $f\in H\cap K$ that meet $e$, $s-2$ of them by \stepref{s414degmeet}, and the $f\in K\setminus H$
  disjoint from $e$. Of the $\binom{m-2}2$ members of $K$ disjoint from $e$
  (\stepref{s414degall}), $h-1-(s-2)$ lie in $H$. The total is
  $(s-2)+\binom{m-2}2-(h-s+1)$.~\qed\end{proof}
 \step{s414degout}{\case{$e\notin H$.}\prove{The count is $2(m-2)+h-2s$.}}
 \begin{proof}\pf\ By Definition~\ref{def:414-switch}, the neighbours are the $f\in K\setminus H$
  with $f\ne e$ that meet $e$, which number $2(m-2)-s$ by \stepref{s414degmeet} and
  \stepref{s414degall}, and the $f\in H\cap K$ disjoint from $e$, which number
  $h-s$.~\qed\end{proof}
 \qedstep
 \begin{proof}\pf\ The two cases are exhaustive.~\qed\end{proof}
\end{proof}

\begin{namedresult}{Lemma}{(constant pair sums)}\label{lem:414-constsum}
 \lean{Finset.two_mul_eq_of_forall_add_eq}
 \uses{}
 \assume{$M$ is a finite set with at least 3 elements, $d:M\to\ints$ and $\sigma\in\ints$, and
 $d(i)+d(j)=\sigma$ for all distinct $i,j\in M$.}
 \prove{$2d(i)=\sigma$ for every $i\in M$.}
\end{namedresult}
\begin{proof}
 \pfsketch\ Compare two pairs with a common member.
 \step{s414csconst}{$d(i)=d(j)$ for all $i,j\in M$.}
 \begin{proof}\pf\ For $i\ne j$ pick $l\in M\setminus\{i,j\}$; then
  $d(i)+d(l)=\sigma=d(j)+d(l)$ by \asm{0}.~\qed\end{proof}
 \qedstep
 \begin{proof}\pf\ For $i\in M$ pick $j\ne i$; then $\sigma=d(i)+d(j)=2d(i)$ by
  \stepref{s414csconst}.~\qed\end{proof}
\end{proof}

\begin{namedresult}{Theorem}{(SE4: regular graphs on 8 vertices)}\label{res:414-SE4}
 \lean{SimpleGraph.IsRegularOfDegree.nonempty_iso_cycle}
 \lean{SimpleGraph.IsRegularOfDegree.eq_bot_of_zero}
 \lean{SimpleGraph.IsRegularOfDegree.nonempty_iso_fourK2}
 \lean{SimpleGraph.IsRegularOfDegree.nonempty_iso_of_two_of_card_eq_eight}
 \lean{SimpleGraph.IsRegularOfDegree.nonempty_iso_of_three_of_card_eq_eight}
 \uses{}
 \assume{$Z$ is a finite graph.}
 \prove{\begin{pfenum}
  \item If $Z$ is connected and 2-regular, then $Z$ is a cycle.
  \item If $Z$ is $d$-regular on 8 vertices with $d\le3$, then $Z$ is isomorphic to one of the
   following eleven graphs: for $d=0$, $\comp{K_8}$; for $d=1$, $4K_2$; for $d=2$, $C_8$,
   $C_3+C_5$ or $2C_4$; for $d=3$, $2K_4$, the cube $Q_3$, the Wagner graph $M_8$, or the graphs
   $Y_1$, $Y_2$, $Y_4$ of \path{cgss/proofs/reg8/reg8.tex}, Definition~1.1.
 \end{pfenum}}
 This is an input, proved in \path{cgss/proofs/reg8/reg8.tex} (Lemma~2.1 and Theorem~1.2); in Lean its
 graphs are \texttt{SimpleGraph.fourK2}, \texttt{cycleThreeFive}, \texttt{twoC4}, \texttt{twoK4},
 \texttt{cubeGraph}, \texttt{wagnerGraph}, \texttt{truncatedK33Graph},
 \texttt{twiceTruncatedK4Graph} and \texttt{twoDiamondsGraph} on \texttt{Fin 8}. Its vertex $i$ is our vertex $i+1$.
\end{namedresult}
\begin{proof}
 \pfsketch\ We cite the structured case split of \path{cgss/proofs/reg8/reg8.tex}.
 \qedstep
 \begin{proof}\pf\ Lemma~2.1 of \path{cgss/proofs/reg8/reg8.tex} for assertion 1, and its Theorem~1.2
  for assertion 2.~\qed\end{proof}
\end{proof}

\begin{namedresult}{Lemma}{(2-regular graphs on 5 and 6 vertices)}\label{lem:414-tworeg}
 \lean{SimpleGraph.IsRegularOfDegree.nonempty_iso_cycle_of_two_of_card_eq_five}
 \lean{SimpleGraph.IsRegularOfDegree.nonempty_iso_of_two_of_card_eq_six}
 \uses{res:414-SE4}
 \assume{$Z$ is a 2-regular graph on $p\in\{5,6\}$ vertices.}
 \prove{If $p=5$ then $Z\cong C_5$; if $p=6$ then $Z\cong C_6$ or $Z\cong C_3+C_3$.}
\end{namedresult}
\begin{proof}
 \pfsketch\ The components are cycles of length at least 3.
 \step{s414twocomp}{Each component of $Z$ is a connected 2-regular graph with at least 3
  vertices.}
 \begin{proof}\pf\ The neighbours of a vertex lie in its component, so the component is
  2-regular; a vertex and its two neighbours are three vertices.~\qed\end{proof}
 \step{s414twocyc}{Each component of $Z$ is a cycle.}
 \begin{proof}\pf\ Theorem~\ref{res:414-SE4}, assertion 1, with
  \stepref{s414twocomp}.~\qed\end{proof}
 \qedstep
 \begin{proof}\pf\ By \stepref{s414twocomp} the orders of the components are at least 3 and add up
  to $p$: for $p=5$ this is $5$ alone, and for $p=6$ it is $6$ or $3+3$. Apply
  \stepref{s414twocyc}.~\qed\end{proof}
\end{proof}

\begin{namedresult}{Computation}{C22}\label{comp:C22}
 \lean{Computations.C22.e6_switch_cycle5}\lean{Computations.C22.e7_switch_cycle6}
 \lean{Computations.C22.e7_switch_twoC3}\lean{Computations.C22.e8_switch_fourK2}
 \lean{Computations.C22.e8_switch_twoC4}\lean{Computations.C22.e8_switch_cube}
 \lean{Computations.C22.e8_switch_twoDiamonds}\lean{Computations.C22.e8_switch_cycle8}
 \lean{Computations.C22.e8_switch_wagner}\lean{Computations.C22.e8_switch_twiceTruncatedK4}
 \lean{Computations.C22.e8_switch_cycleThreeFive}\lean{Computations.C22.e8_switch_truncatedK33}
 \lean{Computations.C22.e8_switch_twoK4}\lean{Computations.C22.e8_switch_empty}
 \lean{Computations.C22.cliqueProfile_t8}\lean{Computations.C22.cliqueProfile_changI}
 \lean{Computations.C22.cliqueProfile_changII}\lean{Computations.C22.cliqueProfile_changIII}
 \lean{Computations.C22.cliqueCount_shrikhande}
 \lean{Computations.C22.cliqueCount_lineGraph_completeBipartite}
 \uses{ex:4.5,ex:4.7,ex:4.9,def:414-switch,res:414-SE4,facts:414}
 \assume{$H_5$ is the edge set of the 5-cycle $1,2,3,4,5$; $H_6$ that of the 6-cycle
 $1,\dots,6$; $H_{3,3}=\{12,13,23,45,46,56\}$; the graphs of Theorem~\ref{res:414-SE4},
 assertion 2, are placed on $\{1,\dots,8\}$; and $\cprof(Z)$ is the multiset of the numbers of
 4-cliques of $Z$ through its vertices.}
 \prove{\begin{pfenum}
  \item $L(\Delta_6)^{H_5}$ is isomorphic to the Petersen graph of Example~\ref{ex:4.9}.
  \item $L(\Delta_7)^{H_6}$ is isomorphic to the Shrikhande graph of Example~\ref{ex:4.7}, and
   $L(\Delta_7)^{H_{3,3}}$ to $L(K_{4,4})$.
  \item For the graphs $Z$ of Theorem~\ref{res:414-SE4}, assertion 2, $L(\Delta_8)^{E(Z)}$ is
   isomorphic to $T(8)$ for $Z=\comp{K_8},2K_4$; to the Chang graph I of Example~\ref{ex:4.5}
   for $Z=4K_2,2C_4,Q_3,Y_4$; to the Chang graph II for $Z=C_8,M_8,Y_2$; and to the Chang graph
   III for $Z=C_3+C_5,Y_1$.
  \item $\cprof(T(8))=\ms{40^{28}}$, and $\cprof$ of the Chang graphs I, II and III is
   $\ms{32^4,36^{24}}$, $\ms{34^{24},36^4}$ and $\ms{33^{10},35^{18}}$.
  \item The Shrikhande graph has no 4-clique, and $L(K_{4,4})$ has 8.
 \end{pfenum}}
 These are SE5 and the part of SE6 that we need (\path{cgss/prerequisites.md}), with the
 4-clique profile in place of the maximal cliques (\path{cgss/certificates/README.md}, C22).
\end{namedresult}
\begin{proof}
 \pfsketch\ Each isomorphism is an explicit permutation of at most 28 vertices, found by
 SageMath (\path{cgss/certificates/c22_switching.sage}, certificates in
 \path{cgss/certificates/data/c22.json}); Lean checks it by \path{Computations.checkGraphIso} and
 \texttt{decide +kernel}. The profiles are evaluated by \path{Computations.cliqueProfile}, which
 enumerates only the 4-subsets that are cliques.
 \step{s414c22iso}{Assertions 1 to 3 hold.}
 \begin{proof}\raggedright\pf\ For each of the fourteen switched graphs, the adjacency table and
  the certificate permutation pass \path{Computations.checkGraphIso}; apply
  \path{Computations.nonempty_graphOfMatrix_iso}. For $Z=\comp{K_8}$ the permutation is the
  identity, since $L(\Delta_8)^{\varnothing}=T(8)$.~\qed\end{proof}
 \step{s414c22prof}{Assertions 4 and 5 hold.}
 \begin{proof}\raggedright\pf\ \path{Computations.cliqueProfile} and \path{Computations.cliqueCount} on the
  adjacency tables, with their soundness theorems; the profiles of assertion 4 are those of
  \path{Computations.Examples.t8_clique_profile} in
  \path{LineGraphs/Computations/CliqueBenchmarkT8.lean} and of its analogues for the Chang graphs in
  \path{LineGraphs/Computations/CliqueBenchmarkChangI.lean},
  \path{LineGraphs/Computations/CliqueBenchmarkChangII.lean} and
  \path{LineGraphs/Computations/CliqueBenchmarkChangIII.lean}.~\qed\end{proof}
 \qedstep
 \begin{proof}\pf\ \stepref{s414c22iso} and \stepref{s414c22prof}.~\qed\end{proof}
\end{proof}

\begin{namedresult}{Lemma}{(regular switchings of $T(5)$)}\label{lem:414-E6}
 \lean{SimpleGraph.nonempty_iso_petersen_of_switch_triangular5}
 \uses{def:414-switch,lem:414-switch,lem:414-degree,lem:414-constsum,lem:414-tworeg,comp:C22,ex:4.9,facts:414}
 \assume{$H$ is a subset of $\Ecal_6$, and the graph $L(\Delta_6)^H$ is 3-regular.}
 \prove{$L(\Delta_6)^H$ is isomorphic to the Petersen graph of Example~\ref{ex:4.9}.}
\end{namedresult}
\begin{proof}
 \pfsketch\ The degree count forces $H$ to be the edge set of a 2-regular graph on
 $\{1,\dots,5\}$, which is a 5-cycle.
 \pflet{$h=|H|$ and $M=\{1,\dots,5\}$.}
 \step{s414esixsum}{$2(\deg_H(i)+\deg_H(j))=h+3$ for all distinct $i,j\in M$.}
 \begin{proof}\pf\ Every vertex of $L(\Delta_6)$ lies in $K=\Ecal_6$. Lemma~\ref{lem:414-degree} with
  $m=5$, $\binom32=3$ and $2(m-2)=6$ gives degree $2s-h$ if $\{i,j\}\in H$ and $6+h-2s$ if not;
  both equal 3 by \asm{0}.~\qed\end{proof}
 \step{s414esixreg}{$4\deg_H(i)=h+3$ for every $i\in M$.}
 \begin{proof}\pf\ Lemma~\ref{lem:414-constsum}, applied to $d=2\deg_H$ on $M$ with $\sigma=h+3$,
  and \stepref{s414esixsum}: $2\cdot2\deg_H(i)=h+3$.~\qed\end{proof}
 \step{s414esixtwo}{$(M,H)$ is 2-regular.}
 \begin{proof}\pf\ Summing \stepref{s414esixreg} over $M$, with $\sum_i\deg_H(i)=2h$
  (\path{SimpleGraph.sum_degrees_eq_twice_card_edges}), gives $8h=5(h+3)$, so $h=5$ and
  $\deg_H(i)=2$.~\qed\end{proof}
 \step{s414esixrelabel}{$L(\Delta_6)^H\cong L(\Delta_6)^{H_5}$.}
 \begin{proof}\pf\ By \stepref{s414esixtwo} and Lemma~\ref{lem:414-tworeg}, some bijection $\sigma$ of
  $M$ carries $H$ onto $H_5$; extended by fixing 6, 7, 8, it is an automorphism of $\Delta_6$. Apply
  Lemma~\ref{lem:414-switch}, assertion 3.~\qed\end{proof}
 \qedstep
 \begin{proof}\pf\ \stepref{s414esixrelabel} and Computation~\ref{comp:C22}, assertion
  1.~\qed\end{proof}
\end{proof}

\begin{namedresult}{Lemma}{(regular switchings of $T(6)$ with an isolated vertex)}\label{lem:414-E7}
 \lean{SimpleGraph.nonempty_iso_shrikhande_or_lineGraph_completeBipartite_of_switch_triangular6AddOne}
 \uses{def:414-switch,lem:414-switch,lem:414-degree,lem:414-constsum,lem:414-tworeg,comp:C22,ex:4.7}
 \assume{$H$ is a subset of $\Ecal_7$, and the graph $L(\Delta_7)^H$ is 6-regular.}
 \prove{$L(\Delta_7)^H$ is isomorphic to the Shrikhande graph of Example~\ref{ex:4.7} or to
 $L(K_{4,4})$.}
\end{namedresult}
\begin{proof}
 \pfsketch\ After complementing $H$ we may assume that $\varepsilon\notin H$; then the degree
 of $\varepsilon$ gives $|H|=6$, and the degree count makes $H$ 2-regular on
 $\{1,\dots,6\}$.
 \pflet{$M=\{1,\dots,6\}$, $K=\Ecal_7\setminus\{\varepsilon\}$, and $H^*=H$ if
 $\varepsilon\notin H$, $H^*=\Ecal_7\setminus H$ if $\varepsilon\in H$.}
 \step{s414esevcompl}{$\varepsilon\notin H^*$, $H^*\sbs K$, and $L(\Delta_7)^H=L(\Delta_7)^{H^*}$.}
 \begin{proof}\pf\ The definition of $H^*$, and Lemma~\ref{lem:414-switch}, assertion 1, when
  $H^*=\Ecal_7\setminus H$.~\qed\end{proof}
 \step{s414esevsix}{$|H^*|=6$.}
 \begin{proof}\pf\ $\varepsilon$ is disjoint from every other member of $\Ecal_7$, so it is
  isolated in $L(\Delta_7)$; by Definition~\ref{def:414-switch} and \stepref{s414esevcompl} its
  neighbours in $L(\Delta_7)^{H^*}$ are the members of $H^*$. Its degree is 6 by \asm{0} and
  \stepref{s414esevcompl}.~\qed\end{proof}
 \step{s414esevsum}{$\deg_{H^*}(i)+\deg_{H^*}(j)=4$ for all distinct $i,j\in M$.}
 \begin{proof}\pf\ Let $e=\{i,j\}$ and $s=\deg_{H^*}(i)+\deg_{H^*}(j)$. By
  Lemma~\ref{lem:414-degree} with $m=6$, $\binom42=6$, $2(m-2)=8$ and $h=6$
  (\stepref{s414esevsix}), $e$ has $2s-3$ neighbours in $K$ if $e\in H^*$ and $14-2s$ if not. By
  Definition~\ref{def:414-switch}, $\varepsilon$ is a further neighbour exactly when $e\in H^*$.
  So the degree is $2s-2$ or $14-2s$, and it is 6.~\qed\end{proof}
 \step{s414esevtwo}{$(M,H^*)$ is 2-regular.}
 \begin{proof}\pf\ Lemma~\ref{lem:414-constsum} with $\sigma=4$ and
  \stepref{s414esevsum}.~\qed\end{proof}
 \step{s414esevrelabel}{$L(\Delta_7)^{H^*}\cong L(\Delta_7)^{H_6}$ or $L(\Delta_7)^{H^*}\cong
  L(\Delta_7)^{H_{3,3}}$.}
 \begin{proof}\pf\ By \stepref{s414esevtwo} and Lemma~\ref{lem:414-tworeg}, some bijection
  $\sigma$ of $M$ carries $H^*$ onto $H_6$ or onto $H_{3,3}$; extended by fixing 7 and 8, it is
  an automorphism of $\Delta_7$. Apply Lemma~\ref{lem:414-switch}, assertion 3.~\qed\end{proof}
 \qedstep
 \begin{proof}\pf\ \stepref{s414esevcompl}, \stepref{s414esevrelabel} and
  Computation~\ref{comp:C22}, assertion 2.~\qed\end{proof}
\end{proof}

\begin{namedresult}{Lemma}{(regular switchings of $T(8)$)}\label{lem:414-E8}
 \lean{SimpleGraph.nonempty_iso_triangular8_or_chang_of_switch_triangular8}
 \uses{def:414-switch,lem:414-switch,lem:414-degree,lem:414-constsum,res:414-SE4,comp:C22,ex:4.5,facts:414}
 \assume{$H$ is a subset of $\Ecal_8$, and the graph $L(\Delta_8)^H$ is 12-regular.}
 \prove{$L(\Delta_8)^H$ is isomorphic to $T(8)$ or to one of the Chang graphs I, II and III of
 Example~\ref{ex:4.5}.}
\end{namedresult}
\begin{proof}
 \pfsketch\ Here the degree count says only that $H$ is regular; after complementing, its valency
 is at most 3, and SE4 leaves eleven graphs.
 \pflet{$h=|H|$ and $M=\{1,\dots,8\}$.}
 \step{s414eeightsum}{$2(\deg_H(i)+\deg_H(j))=h$ for all distinct $i,j\in M$.}
 \begin{proof}\pf\ Every vertex of $L(\Delta_8)$ lies in $K=\Ecal_8$. Lemma~\ref{lem:414-degree} with
  $m=8$, $\binom62=15$ and $2(m-2)=12$ gives degree $2s-h+12$ if $\{i,j\}\in H$ and $12+h-2s$ if
  not; both equal 12 by \asm{0}.~\qed\end{proof}
 \step{s414eeightreg}{$(M,H)$ is regular of some valency $d\le7$.}
 \begin{proof}\pf\ Lemma~\ref{lem:414-constsum}, applied to $d'=2\deg_H$ on $M$ with $\sigma=h$, and
  \stepref{s414eeightsum} give $4\deg_H(i)=h$ for every $i\in M$; so $d=h/4$, and $d\le7$ since
  $|M|=8$.~\qed\end{proof}
 \pflet{$H^*=H$ if $d\le3$, and $H^*=\Ecal_8\setminus H$ if $d\ge4$.}
 \step{s414eeightcompl}{$(M,H^*)$ is regular of valency at most 3, and
  $L(\Delta_8)^H=L(\Delta_8)^{H^*}$.}
 \begin{proof}\pf\ $\Ecal_8\setminus H$ is the edge set of the complement of $(M,H)$, which is
  $(7-d)$-regular (\path{SimpleGraph.IsRegularOfDegree.compl}); apply
  Lemma~\ref{lem:414-switch}, assertion 1.~\qed\end{proof}
 \step{s414eeightrelabel}{$L(\Delta_8)^{H^*}\cong L(\Delta_8)^{E(Z)}$ for one of the eleven graphs $Z$ of
  Theorem~\ref{res:414-SE4}, assertion 2, placed on $M$.}
 \begin{proof}\pf\ By \stepref{s414eeightcompl} and Theorem~\ref{res:414-SE4}, some bijection
  $\sigma$ of $M$ is an isomorphism from $(M,H^*)$ to such a $Z$; it is an automorphism of
  $\Delta_8=K_8$. Apply Lemma~\ref{lem:414-switch}, assertion 3.~\qed\end{proof}
 \qedstep
 \begin{proof}\pf\ \stepref{s414eeightcompl}, \stepref{s414eeightrelabel} and
  Computation~\ref{comp:C22}, assertion 3.~\qed\end{proof}
\end{proof}

\subsubsection*{The Schl\"afli and Clebsch graphs}

\begin{namedresult}{Lemma}{(the triangle property from the parameters, GQ1)}\label{lem:414-GQ1}
 \lean{SimpleGraph.IsSRGWith.hasTriangleProperty}
 \uses{def:triangle-property}
 \assume{$Z$ is strongly regular with parameters $(n,k,1,c)$, and $n+3=3k$.}
 \prove{$Z$ has the triangle property.}
 \new\ This is the converse half of GQ1 of \path{cgss/prerequisites.md}, with the triangle
 property of Section 2 in place of a generalized quadrangle whose lines are the triangles;
 [BCN] Lemma 1.15.1 says only ``Straightforward''. The hypothesis is $n-3=3(k-2)$, written
 without subtraction as in Lean.
\end{namedresult}
\begin{proof}
 \pfsketch\ The edges leaving a triangle reach every other vertex exactly once.
 \step{s414gqone}{Every edge $xy$ has exactly one common neighbour $z$.}
 \begin{proof}\pf\ The parameter $a=1$.~\qed\end{proof}
 \step{s414gqmost}{A vertex $v\notin\{x,y,z\}$ is adjacent to at most one of $x,y,z$.}
 \begin{proof}\pf\ If $v$ were adjacent to two of them, say to $x$ and $y$, then $v$ and $z$ would
  be two common neighbours of the edge $xy$, against \stepref{s414gqone}; the other pairs are
  edges of the triangle as well.~\qed\end{proof}
 \qedstep
 \begin{proof}\pf\ Each of $x,y,z$ has $k-2$ neighbours outside $\{x,y,z\}$, so there are
  $3(k-2)=n-3$ edges from the triangle to the other $n-3$ vertices, by \asm{0}; by
  \stepref{s414gqmost} no vertex receives two, so every vertex receives exactly one. With
  \stepref{s414gqone}, this is Definition~\ref{def:triangle-property}.~\qed\end{proof}
\end{proof}

\begin{namedresult}{Definition}{(the folded 5-cube)}\label{def:414-folded}
 \lean{SimpleGraph.foldedFiveCube}\lean{SimpleGraph.compl_foldedFiveCube_adj_iff}
 \uses{}
 \assume{$\Omega_5=\{1,\dots,5\}$, and $V_{16}$ is the set of the subsets of $\Omega_5$ of even
 size.}
 \pflet{The \emph{folded 5-cube} $\square_5$ is the graph on $V_{16}$ in which $s\sim t$ when
 $|s\sd t|=4$.}
 This is Definition 1.1 of \path{cgss/proofs/clebsch/clebsch.tex}. $V_{16}$ consists of $\varnothing$,
 the ten 2-subsets and the five 4-subsets of $\Omega_5$. In Lean the vertex type is
 \texttt{$\{$s : Finset (Fin 5) // Even s.card$\}$}, whose \texttt{i} is our $i+1$. In the complement, distinct $s,t$ are adjacent
 exactly when $|s\sd t|=2$, since $|s\sd t|$ is even and at most 5.
\end{namedresult}

\begin{namedresult}{Definition}{(the double-six graph)}\label{def:414-gq24}
 \lean{SimpleGraph.gq24}\lean{SimpleGraph.compl_gq24_adj_iff}
 \uses{}
 \assume{$\Omega_6=\{1,\dots,6\}$, and $V_{27}$ is the set of the 27 symbols $a_i$, $b_i$
 ($i\in\Omega_6$) and $c_{ij}=c_{ji}$ ($i\ne j$ in $\Omega_6$).}
 \pflet{$Q_{27}$ be the graph on $V_{27}$ in which $a_i\sim b_j$ for $i\ne j$; $a_i\sim c_{jl}$ and
 $b_i\sim c_{jl}$ for $i\in\{j,l\}$; $c_{ij}\sim c_{lm}$ when $\{i,j\}\cap\{l,m\}=\varnothing$;
 and there are no other edges.}
 This is the graph $Q$ of \path{cgss/proofs/clebsch/clebsch.tex}, Definition 1.2, the collinearity
 graph of the generalized quadrangle of order $(2,4)$; we write $Q_{27}$, since $Q$ is $2I+A$ in
 this section. In Lean the vertex type is \texttt{Fin 2} $\times$ \texttt{Fin 6} $\oplus$
 \texttt{T} with \texttt{T = $\{$s : Finset (Fin 6) // s.card = 2$\}$}; $(0,i)$ is $a_{i+1}$, $(1,i)$
 is $b_{i+1}$, and $\{i,j\}$ is $c_{i+1,j+1}$. The second Lean name spells out adjacency in
 $\comp{Q_{27}}$ by the same cases.
\end{namedresult}

\begin{namedresult}{Computation}{(the models of SE2 and SE3)}\label{comp:414-models}
 \lean{SimpleGraph.foldedFiveCube_isSRGWith}\lean{SimpleGraph.gq24_isSRGWith}
 \uses{def:414-folded,def:414-gq24,comp:C9,ex:4.6,ex:4.8}
 \assume{$X_{\mathrm{Cl}}$ and $X_{\mathrm{Sch}}$ are the Clebsch and Schl\"afli graphs of
 Examples~\ref{ex:4.8} and~\ref{ex:4.6}, that is, the graphs of the vectors of their
 representations (Computation~\ref{comp:C9}).}
 \prove{\begin{pfenum}
  \item $\square_5$ is strongly regular with parameters $(16,5,0,2)$, and $Q_{27}$ with
   parameters $(27,10,1,5)$.
  \item $X_{\mathrm{Cl}}\cong\comp{\square_5}$ and $X_{\mathrm{Sch}}\cong\comp{Q_{27}}$.
 \end{pfenum}}
 The Clebsch and Schl\"afli graphs stay the graphs of the vectors; $\square_5$ and $Q_{27}$ are
 the combinatorial models of \path{cgss/proofs/clebsch/clebsch.tex}. The isomorphisms of assertion 2
 are the certificates of Computation~\ref{comp:C9}, which carry their Lean names there; we
 repeat them here beside the parameters.
\end{namedresult}
\begin{proof}
 \pfsketch\ These are finite checks on at most 27 vertices, as for Computation C9.
 \step{s414modpar}{Assertion 1 holds.}
 \begin{proof}\pf\ In Lean, \texttt{decide} counts the neighbours of every vertex and the common
  neighbours of every pair, at most $27^3$ triples. By hand, Lemmas 3.1 and 4.2 of
  \path{cgss/proofs/clebsch/clebsch.tex} count them at one vertex of each orbit of an automorphism
  group.~\qed\end{proof}
 \step{s414modiso}{Assertion 2 holds.}
 \begin{proof}\pf\ Computation~\ref{comp:C9}, its last sentence.~\qed\end{proof}
 \qedstep
 \begin{proof}\pf\ \stepref{s414modpar} and \stepref{s414modiso}.~\qed\end{proof}
\end{proof}

\begin{namedresult}{Lemma}{(the double-six graph from its parameters)}\label{lem:414-gq24iso}
 \lean{SimpleGraph.IsSRGWith.nonempty_iso_gq24}
 \uses{lem:414-GQ1,lem:triangle-iso,def:414-gq24,comp:414-models}
 \assume{$Z$ is a strongly regular graph with parameters $(27,10,1,5)$.}
 \prove{$Z\cong Q_{27}$.}
 \new\ SE3 of \path{cgss/prerequisites.md} identifies $Z$ by GQ2, the classification of the
 generalized quadrangles of order $(2,4)$; we use instead the isomorphism from the order of
 Section 2, as in Theorem 1.4 of \path{cgss/proofs/clebsch/clebsch.tex}.
\end{namedresult}
\begin{proof}
 \pfsketch\ Both graphs have the triangle property, and Section 2 shows that such a graph is
 determined by its order.
 \step{s414gqtri}{$Z$ and $Q_{27}$ have the triangle property.}
 \begin{proof}\pf\ Lemma~\ref{lem:414-GQ1}, with \asm{0} and Computation~\ref{comp:414-models},
  assertion 1, since $27+3=3\cdot10$.~\qed\end{proof}
 \step{s414gqedge}{$Z$ and $Q_{27}$ have 27 vertices each, an edge, and no vertex adjacent to
  every other vertex.}
 \begin{proof}\pf\ Both are 10-regular on 27 vertices, by \asm{0} and
  Computation~\ref{comp:414-models}, assertion 1, and $0<10<26$.~\qed\end{proof}
 \qedstep
 \begin{proof}\pf\ Lemma~\ref{lem:triangle-iso}, with \stepref{s414gqtri} and
  \stepref{s414gqedge}.~\qed\end{proof}
\end{proof}

\begin{namedresult}{Lemma}{(SE3: the Schl\"afli graph)}\label{lem:414-SE3}
 \lean{SimpleGraph.IsSRGWith.nonempty_iso_schlafli}
 \uses{lem:414-iso,lem:414-gq24iso,comp:414-models,ex:4.6,facts:414}
 \assume{$X$ is a strongly regular graph with parameters $(27,16,10,8)$.}
 \prove{$X$ is isomorphic to the Schl\"afli graph of Example~\ref{ex:4.6}.}
 \new\ SE3 of \path{cgss/prerequisites.md} goes through GQ1 and GQ2; we go through
 Lemma~\ref{lem:414-gq24iso} and the certificate of Computation~\ref{comp:414-models}.
\end{namedresult}
\begin{proof}
 \pfsketch\ We will pass to the complement, which is the double-six graph.
 \pflet{$X_{\mathrm{Sch}}$ be the Schl\"afli graph of Example~\ref{ex:4.6}.}
 \step{s414sthreecompl}{$\comp X$ is strongly regular with parameters $(27,10,1,5)$.}
 \begin{proof}\pf\ \path{SimpleGraph.IsSRGWith.compl}: $(27,27-16-1,27-(32-8)-2,27-(32-10))$, with no
  truncated subtraction.~\qed\end{proof}
 \step{s414sthreeiso}{$\comp X\cong Q_{27}$.}
 \begin{proof}\pf\ Lemma~\ref{lem:414-gq24iso} with \stepref{s414sthreecompl}.~\qed\end{proof}
 \qedstep
 \begin{proof}\pf\ Lemma~\ref{lem:414-iso}, assertion 1 (\texttt{SimpleGraph.Iso.compl}), applied
  to \stepref{s414sthreeiso}, with \path{compl_compl}, gives $X\cong\comp{Q_{27}}$; and
  $\comp{Q_{27}}\cong X_{\mathrm{Sch}}$ by Computation~\ref{comp:414-models},
  assertion 2.~\qed\end{proof}
\end{proof}

\begin{namedresult}{Lemma}{(the folded 5-cube from its parameters)}\label{lem:414-foldiso}
 \lean{SimpleGraph.IsSRGWith.nonempty_iso_foldedFiveCube}
 \uses{def:414-folded}
 \assume{$Z$ is a strongly regular graph with parameters $(16,5,0,2)$.}
 \prove{$Z\cong\square_5$.}
 This is the first half of SE2 of \path{cgss/prerequisites.md}, after [C17] Theorem 2.2, and
 Theorem 1.3 of \path{cgss/proofs/clebsch/clebsch.tex}. [BCN] 3.11 (iii) prints these parameters as
 $(16,5,0,1)$; they are $(16,5,0,2)$.
\end{namedresult}
\begin{proof}
 \pfsketch\ The non-neighbours of a vertex $u$ correspond to the 2-subsets of its
 neighbourhood, and adjacency among them is disjointness.
 \pflet{$u$ be a vertex of $Z$, $N$ its neighbourhood and $O=V(Z)\setminus(\{u\}\cup N)$.}
 \step{s414foldcocl}{$N$ is a coclique of 5 vertices, and $|O|=10$.}
 \begin{proof}\pf\ $|N|=5$ and $|O|=16-1-5$. Two adjacent vertices of $N$ would have the common
  neighbour $u$, against $a=0$.~\qed\end{proof}
 \step{s414foldphi}{For $w\in O$, $\phi(w)=N(w)\cap N$ has exactly two elements.}
 \begin{proof}\pf\ $w\ne u$ and $w\not\sim u$, so they have $c=2$ common neighbours, which are
  the members of $N(w)\cap N$.~\qed\end{proof}
 \step{s414foldbij}{$\phi$ is a bijection from $O$ onto $\binom N2$.}
 \begin{proof}\pf\ If $w\ne w'$ in $O$ had $\phi(w)=\phi(w')=\{x,y\}$, the distinct non-adjacent
  vertices $x,y$ (\stepref{s414foldcocl}) would have the three common neighbours $u,w,w'$,
  against $c=2$. So $\phi$ is injective, and $|O|=10=\binom52$.~\qed\end{proof}
 \step{s414foldmeet}{\assume{$w,w'\in O$ and $\phi(w)\cap\phi(w')\ne\varnothing$.}
  \prove{$w\not\sim w'$.}}
 \begin{proof}\pf\ Take $x\in\phi(w)\cap\phi(w')$ by \asm{1}; an edge $ww'$ would have the common
  neighbour $x$, against $a=0$.~\qed\end{proof}
 \step{s414foldO}{Each $w\in O$ has exactly 3 neighbours in $O$.}
 \begin{proof}\pf\ $w$ has 5 neighbours; two lie in $N$ (\stepref{s414foldphi}), and $u$ is not
  one.~\qed\end{proof}
 \step{s414folddisj}{Distinct $w,w'\in O$ are adjacent if and only if $\phi(w)$ and $\phi(w')$
  are disjoint.}
 \begin{proof}\pf\ Exactly $\binom32=3$ members of $\binom N2$ are disjoint from $\phi(w)$, so by
  \stepref{s414foldbij} exactly 3 vertices $w'\in O$ have $\phi(w')$ disjoint from $\phi(w)$. By
  \stepref{s414foldmeet} the 3 neighbours of $w$ in $O$ (\stepref{s414foldO}) are among them;
  so they are all of them.~\qed\end{proof}
 \pflet{$\beta:N\to\Omega_5$ be a bijection, and $f:V(Z)\to V_{16}$ the map with $f(u)=\varnothing$,
 $f(x)=\Omega_5\setminus\{\beta(x)\}$ for $x\in N$ and $f(w)=\beta(\phi(w))$ for $w\in O$.}
 \step{s414foldf}{$f$ is a bijection.}
 \begin{proof}\pf\ $f$ maps $u$ to $\varnothing$, $N$ onto the five 4-subsets, since $\beta$ is a
  bijection, and $O$ onto the ten 2-subsets, by \stepref{s414foldbij}; these are all of $V_{16}$
  (Definition~\ref{def:414-folded}).~\qed\end{proof}
 \step{s414foldsd}{For $i\in\Omega_5$ and 2-subsets $P\ne P'$ of $\Omega_5$,
  $|(\Omega_5\setminus\{i\})\sd P|$ is 4 if $i\in P$ and 2 if $i\notin P$, and
  $|P\sd P'|=4-2|P\cap P'|$.}
 \begin{proof}\pf\ $|s\sd s'|=|s|+|s'|-2|s\cap s'|$, and $(\Omega_5\setminus\{i\})\cap P$ has 1
  element if $i\in P$ and 2 if not.~\qed\end{proof}
 \step{s414foldadj}{Distinct vertices $p,q$ of $Z$ are adjacent if and only if
  $|f(p)\sd f(q)|=4$.}
 \begin{proof}
  \step{s414foldcaseone}{\case{$p=u$ and $q=x\in N$.}\prove{Both hold.}}
  \begin{proof}\pf\ $u\sim x$, and $|\varnothing\sd(\Omega_5\setminus\{\beta(x)\})|=4$.~\qed\end{proof}
  \step{s414foldcasetwo}{\case{$p=u$ and $q=w\in O$.}\prove{Neither holds.}}
  \begin{proof}\pf\ $w\notin N$, and $|\varnothing\sd f(w)|=|f(w)|=2$.~\qed\end{proof}
  \step{s414foldcasethree}{\case{$p=x$ and $q=y$ lie in $N$.}\prove{Neither holds.}}
  \begin{proof}\pf\ $x\not\sim y$ by \stepref{s414foldcocl}, and
   $|f(x)\sd f(y)|=|\{\beta(x),\beta(y)\}|=2$.~\qed\end{proof}
  \step{s414foldcasefour}{\case{$p=x\in N$ and $q=w\in O$.}\prove{The two conditions agree.}}
  \begin{proof}\pf\ $x\sim w$ exactly when $x\in\phi(w)$, that is, when $\beta(x)\in f(w)$, and
   by \stepref{s414foldsd} this is when $|f(x)\sd f(w)|=4$.~\qed\end{proof}
  \step{s414foldcasefive}{\case{$p=w$ and $q=w'$ lie in $O$.}\prove{The two conditions
   agree.}}
  \begin{proof}\pf\ $w\sim w'$ exactly when $f(w)\cap f(w')=\varnothing$, by
   \stepref{s414folddisj} and the injectivity of $\beta$, and by \stepref{s414foldsd} this is
   when $|f(w)\sd f(w')|=4$.~\qed\end{proof}
  \qedstep
  \begin{proof}\pf\ $V(Z)$ is the disjoint union of $\{u\}$, $N$ and $O$, so up to the order of
   $p$ and $q$ the five cases are exhaustive.~\qed\end{proof}
 \end{proof}
 \qedstep
 \begin{proof}\pf\ By \stepref{s414foldf} and \stepref{s414foldadj}, $f$ is an isomorphism from
  $Z$ to $\square_5$ (Definition~\ref{def:414-folded}).~\qed\end{proof}
\end{proof}

\begin{namedresult}{Lemma}{(SE2: the Clebsch graph)}\label{lem:414-SE2}
 \lean{SimpleGraph.IsSRGWith.nonempty_iso_clebsch}
 \uses{lem:414-iso,lem:414-foldiso,comp:414-models,ex:4.8,facts:414}
 \assume{$X$ is a strongly regular graph with parameters $(16,10,6,6)$.}
 \prove{$X$ is isomorphic to the Clebsch graph of Example~\ref{ex:4.8}.}
 This is SE2 of \path{cgss/prerequisites.md}.
\end{namedresult}
\begin{proof}
 \pfsketch\ We will pass to the complement, which is the folded 5-cube.
 \pflet{$X_{\mathrm{Cl}}$ be the Clebsch graph of Example~\ref{ex:4.8}.}
 \step{s414stwocompl}{$\comp X$ is strongly regular with parameters $(16,5,0,2)$.}
 \begin{proof}\pf\ \path{SimpleGraph.IsSRGWith.compl}: $(16,16-10-1,16-(20-6)-2,16-(20-6))$, with no
  truncated subtraction.~\qed\end{proof}
 \step{s414stwofold}{$\comp X\cong\square_5$.}
 \begin{proof}\pf\ Lemma~\ref{lem:414-foldiso} with \stepref{s414stwocompl}.~\qed\end{proof}
 \qedstep
 \begin{proof}\pf\ Lemma~\ref{lem:414-iso}, assertion 1 (\texttt{SimpleGraph.Iso.compl}), applied
  to \stepref{s414stwofold}, with \path{compl_compl}, gives $X\cong\comp{\square_5}$; and
  $\comp{\square_5}\cong X_{\mathrm{Cl}}$ by Computation~\ref{comp:414-models},
  assertion 2.~\qed\end{proof}
\end{proof}

\subsubsection*{The seven graphs}

\begin{namedresult}{Proposition}{(the seven graphs)}\label{prop:414-seven}
 \lean{SimpleGraph.seidel_seven_properties}\lean{SimpleGraph.seidel_seven_pairwise_not_iso}
 \uses{def:CP,def:represented,lem:4-gram,comp:C9,comp:C22,ex:4.5,ex:4.6,ex:4.7,ex:4.8,ex:4.9,facts:414}
 \assume{$\cS$ is the list of the seven graphs of Examples 4.5 to 4.9: the Chang graphs I, II
 and III, the Schl\"afli graph, the Shrikhande graph, the Clebsch graph and the Petersen
 graph.}
 \prove{\begin{pfenum}
  \item Each $Z\in\cS$ is connected and strongly regular, has $\lambda_{\min}(Z)=-2$, and is
   represented in $\typ E_8$.
  \item No $Z\in\cS$ is isomorphic to $L(K_m)$, to $L(K_{m,m})$ or to $CP(m)$, for any $m$.
  \item No two graphs of $\cS$ are isomorphic.
 \end{pfenum}}
 \repair\ ``The seven graphs'' presupposes that the three Chang graphs are pairwise
 non-isomorphic, which the paper does not show (errata, Gap 12). We will separate them by their
 4-clique profiles (Computation C22), which the kernel can evaluate by extending cliques one
 vertex at a time; their numbers of maximal cliques, used in SE6 of
 \path{cgss/prerequisites.md}, can be enumerated in the same way, but every clique found must
 then be checked for maximality, which costs the kernel much more.
\end{namedresult}
\begin{proof}
 \pfsketch\ The parameters separate all pairs except the ones sharing $(n,k)$, and C22 separates
 those.
 \step{s414sevparams}{Each $Z\in\cS$ is represented in $\typ E_8$, and is strongly regular with
  parameters $(28,12,6,4)$ (Chang), $(27,16,10,8)$, $(16,6,2,2)$, $(16,10,6,6)$ and $(10,3,0,1)$.}
 \begin{proof}\pf\ Examples~\ref{ex:4.5} to~\ref{ex:4.9}, and
  Computation~\ref{comp:C9}.~\qed\end{proof}
 \step{s414sevconn}{Each $Z\in\cS$ is connected.}
 \begin{proof}\pf\ In each case $c\ge1$ by \stepref{s414sevparams}, so any two distinct
  non-adjacent vertices have a common neighbour.~\qed\end{proof}
 \step{s414sevlmin}{$\lambda_{\min}(Z)=-2$ for each $Z\in\cS$.}
 \begin{proof}\pf\ $Z$ is represented in $\Phi(\typ E_8)\sbs\re^8$ (\stepref{s414sevparams}), a
  space of dimension $8<10\le|V(Z)|$; apply Lemma~\ref{lem:4-gram}, assertion
  2.~\qed\end{proof}
 \step{s414sevpar}{$L(K_m)$ has $\binom m2$ vertices and valency $2(m-2)$, $L(K_{m,m})$ has $m^2$
  vertices and valency $2(m-1)$, and $CP(m)$ has $2m$ vertices and valency $2m-2$.}
 \begin{proof}\pf\ An edge $\{i,j\}$ of $K_m$ meets the $2(m-2)$ edges $\{i,l\}$, $\{j,l\}$; an
  edge of $K_{m,m}$ meets $m-1$ others at each end; and Definition~\ref{def:CP}.~\qed\end{proof}
 \step{s414sevnum}{The Schl\"afli, Clebsch and Petersen graphs are isomorphic to none of $L(K_m)$,
  $L(K_{m,m})$, $CP(m)$, and the Chang and Shrikhande graphs to no $CP(m)$.}
 \begin{proof}\pf\ Compare \stepref{s414sevparams} with \stepref{s414sevpar}. For 27 vertices,
  $\binom m2=27$, $m^2=27$ and $2m=27$ have no solution. For 16 vertices, $\binom m2=16$ has none,
  and $m^2=16$, $2m=16$ give valencies 6 and 14, not 10. For 10 vertices, $\binom m2=10$ gives
  valency 6, $m^2=10$ has no solution and $2m=10$ gives 8, not 3. Finally $2m=28$ and $2m=16$
  give valencies 26 and 14, not 12 and 6.~\qed\end{proof}
 \step{s414sevchang}{The Chang graphs I, II and III and $T(8)$ are pairwise non-isomorphic.}
 \begin{proof}\pf\ Their 4-clique profiles differ (Computation~\ref{comp:C22}, assertion 4), and
  an isomorphism preserves the profile (\path{SimpleGraph.Iso.cliqueProfile_eq}).~\qed\end{proof}
 \step{s414sevshr}{The Shrikhande graph is not isomorphic to $L(K_{4,4})$.}
 \begin{proof}\pf\ By Computation~\ref{comp:C22}, assertion 5, $L(K_{4,4})$ has a 4-clique; an
  isomorphism from $L(K_{4,4})$ to the Shrikhande graph would map it to a 4-clique
  (\path{SimpleGraph.IsNClique.map}), which the Shrikhande graph does not have.~\qed\end{proof}
 \step{s414sevlk}{No Chang graph is isomorphic to $L(K_{m,m})$ or to $L(K_m)$ with $m\ne8$, and the
  Shrikhande graph is isomorphic to no $L(K_m)$ and to no $L(K_{m,m})$ with $m\ne4$.}
 \begin{proof}\pf\ By \stepref{s414sevparams} and \stepref{s414sevpar}: $m^2=28$ has no solution,
  $\binom m2=28$ only $m=8$, $\binom m2=16$ none, and $m^2=16$ only $m=4$.~\qed\end{proof}
 \qedstep
 \begin{proof}\pf\ Assertion 1 is \stepref{s414sevparams}, \stepref{s414sevconn} and
  \stepref{s414sevlmin}; assertion 2 is \stepref{s414sevnum}, \stepref{s414sevlk},
  \stepref{s414sevchang} and \stepref{s414sevshr}. For assertion 3, an isomorphism preserves
  $(n,k)$, and by \stepref{s414sevparams} only the Chang graphs share it; apply
  \stepref{s414sevchang}.~\qed\end{proof}
\end{proof}

\begin{namedresult}{Corollary}{(the seven graphs are not line graphs)}\label{cor:414-notline}
 \lean{SimpleGraph.seidelSeven_not_iso_lineGraph}
 \uses{prop:414-seven,lem:414-linegraph}
 \assume{$\cS$ is the list of the seven graphs of Proposition~\ref{prop:414-seven}, and $Y$ is a
 graph, possibly infinite or with isolated vertices.}
 \prove{No $Z\in\cS$ is isomorphic to $L(Y)$.}
 This is the claim of the paper, before Example 4.5, that these strongly regular graphs are not
 line graphs.
\end{namedresult}
\begin{proof}
 \pfsketch\ We will apply Lemma~\ref{lem:414-linegraph}; assertion 2 of
 Proposition~\ref{prop:414-seven} excludes both of its alternatives.
 \step{s414nlcase}{\assume{$Z\in\cS$ and $Z\cong L(Y)$.}\prove{A contradiction.}}
 \begin{proof}
  \step{s414nldet}{$Z$ is connected and strongly regular, and $\det(2I+A(Z))=0$.}
  \begin{proof}\pf\ Proposition~\ref{prop:414-seven}, assertion 1, with $Z\in\cS$ by \asm{1};
   since $\lambda_{\min}(Z)=-2$ is an eigenvalue of $A(Z)$, the matrix $2I+A(Z)$ is
   singular.~\qed\end{proof}
  \step{s414nllg}{$Z\cong L(K_m)$ or $Z\cong L(K_{m,m})$ for some $m\ge2$.}
  \begin{proof}\pf\ Lemma~\ref{lem:414-linegraph}, with \stepref{s414nldet} and
   \asm{1}.~\qed\end{proof}
  \qedstep
  \begin{proof}\pf\ Proposition~\ref{prop:414-seven}, assertion 2, excludes both alternatives
   of \stepref{s414nllg}.~\qed\end{proof}
 \end{proof}
 \qedstep
 \begin{proof}\pf\ Step \stepref{s414nlcase} applies to every $Z\in\cS$.~\qed\end{proof}
\end{proof}

\subsubsection*{The classification}

\begin{namedresult}{Remark}{4.14, restated}\label{rem:414-proof}
 \lean{SimpleGraph.seidel_seven_classification}
 \uses{rem:4.14,def:represented,def:CP,def:3.2,def:3.3,lem:4-gramfactor,lem:4-transport,lem:4-regglg,thm:4.1,thm:4.2,thm:4.4,thm:3.5,lem:sec3-transport,lem:sec3-spans,lem:414-CP,lem:414-srg,comp:C18,lem:414-linegraph,lem:414-lattice,lem:414-apex,lem:414-dual,comp:C21,comp:414-Nc,lem:414-transversal,lem:414-E6,lem:414-E7,lem:414-E8,lem:414-SE2,lem:414-SE3,comp:414-models,prop:414-seven,facts:414}
 \assume{\begin{pfenum}
  \item $X$ is a connected graph, strongly regular with parameters $(n,k,a,c)$.
  \item $\lambda_{\min}(X)=-2$.
  \item $X$ is represented in $\typ E_8$ by $(\Zx_u)$.
  \item For no $m$ is $X$ isomorphic to $L(K_m)$, to $L(K_{m,m})$ or to $CP(m)$.
 \end{pfenum}}
 \prove{\begin{pfenum}
  \item $X$ is isomorphic to exactly one of the Chang graphs I, II and III, the Schl\"afli graph,
   the Shrikhande graph, the Clebsch graph and the Petersen graph.
  \item These seven graphs satisfy hypotheses 1 to 4.
 \end{pfenum}}
 \new\ The paper cites Seidel [19] and gives no proof. We follow the route of
 \path{cgss/prerequisites.md}, Section 7.3 (decision 3 of its Section 13); the departures are
 tagged in the lemmas above. \repair\ The proof fills Gap 12 (errata): it proves the
 classification, and ``exactly one'' is the non-isomorphism of Proposition~\ref{prop:414-seven}.
 In Lean, the statement \texttt{seidel\_seven} of Remark~\ref{rem:4.14} is the
 conjunction of \texttt{seidel\_seven\_classification}, of
 \texttt{seidel\_seven\_pairwise\_not\_iso} and of \texttt{seidel\_seven\_properties}
 (Proposition~\ref{prop:414-seven}), all four in the namespace \texttt{SimpleGraph}; Lean
 names are unique, so the two results carry different ones.
\end{namedresult}
\begin{proof}
 \pfsketch\ We will pin down the parameters, find the type of the lattice roots $R$, exclude the
 types $\typ A$ and $\typ D$ by the line graph theorems, exclude two parameter sets by the
 integrality of the apex, and identify the remaining five cases by switching and by uniqueness.
 \pflet{$A=A(X)$, $Q=2I+A$, $\rho=\rk Q$, and let $\Lambda$, $W$, $R$ be as in
 Lemma~\ref{lem:414-lattice} and $\Zw$ as in Lemma~\ref{lem:414-apex} for $(\Zx_u)$.}
 \step{s414mdet}{$n\ge1$ and $\det Q=0$.}
 \begin{proof}\pf\ A connected graph has a vertex; \asm[2]{0} and
  \path{SimpleGraph.leastEigenvalue_eq_neg_two_iff}.~\qed\end{proof}
 \step{s414mbounds}{$n\le28$ and $k\le16$.}
 \begin{proof}\pf\ Theorem~\ref{thm:4.4}, with \asm[1]{0}, \asm[3]{0} and
  \stepref{s414mdet}.~\qed\end{proof}
 \step{s414mrank}{$\rho=\dim W\le8$.}
 \begin{proof}\pf\ $Q$ is the Gram matrix of $(\Zx_u)$, so $\rho=\dim W$ by
  Lemma~\ref{lem:4-gramfactor}, assertion 1; and $W\sbs\re^8$
  (\path{Submodule.finrank_le}).~\qed\end{proof}
 \step{s414msrg}{$1\le k\le n-2$, $a\le k-1$, $1\le c\le k$, and conditions 2 to 4 of
  Computation~\ref{comp:C18} hold.}
 \begin{proof}\pf\ Lemma~\ref{lem:414-srg} with \asm[1]{0} and \stepref{s414mdet}.~\qed\end{proof}
 \step{s414mparam}{$k(k-a-1)=c(n-k-1)$, condition 1 of Computation~\ref{comp:C18}.}
 \begin{proof}\pf\ \path{SimpleGraph.IsSRGWith.param_eq}, which needs $0<n$ (\stepref{s414mdet});
  it is stated in $\mathbb N$ with truncated subtraction, and by \stepref{s414msrg} neither
  $k-a-1$ nor $n-k-1$ is truncated.~\qed\end{proof}
 \step{s414mCeighteen}{$(n,k,a,c,\rho)$ is one of the 15 tuples of Computation C18.}
 \begin{proof}\pf\ Computation~\ref{comp:C18}: its hypotheses are \stepref{s414msrg},
  \stepref{s414mparam}, and $n\le28\le36$, $k\le16$ and $\rho\le8$ from \stepref{s414mbounds} and
  \stepref{s414mrank}.~\qed\end{proof}
 \step{s414mnotCP}{$(n,k,a,c)\ne(2m,2m-2,2m-4,2m-2)$ for every $m$.}
 \begin{proof}\pf\ Otherwise $X$ is $(2m-2)$-regular on $2m$ vertices, and $X\cong CP(m)$ by
  Lemma~\ref{lem:414-CP}, against \asm[4]{0}.~\qed\end{proof}
 \step{s414mR}{$R$ is a nonempty, indecomposable, star-closed line system spanning $W$, and $X$
  is represented in $R$ by $(\Zx_u)$.}
 \begin{proof}\pf\ Lemma~\ref{lem:414-lattice}, with \asm[1]{0} and \asm[3]{0}.~\qed\end{proof}
 \step{s414mtype}{\pick\ a type $\tau$ among $\typ A_\ell$ ($\ell\ge1$), $\typ D_\ell$
  ($\ell\ge4$), $\typ E_6$, $\typ E_7$, $\typ E_8$, and a bijection $f:\Phi(R)\to\Phi(\tau)$
  preserving inner products.}
 \begin{proof}\pf\ Theorem~\ref{thm:3.5} with \stepref{s414mR}.~\qed\end{proof}
 \step{s414mspan}{\pick\ a linear isometry $F$ from $W$ onto $\opk{span}\Phi(\tau)$ that extends
  $f$.}
 \begin{proof}\pf\ Lemma~\ref{lem:sec3-transport} with \stepref{s414mtype}, since $\Phi(R)$ spans
  $W$ (\stepref{s414mR}).~\qed\end{proof}
 \step{s414mrk}{The rank of $\tau$ is $\rho$.}
 \begin{proof}\pf\ $F$ of \stepref{s414mspan} is an isometry onto $\opk{span}\Phi(\tau)$, and
  $\dim W=\rho$ by \stepref{s414mrank}.~\qed\end{proof}
 \step{s414mAD}{\assume{$\tau$ is $\typ A_\ell$ or $\typ D_\ell$.}\prove{A contradiction.}}
 \begin{proof}
  \step{s414mA}{\case{$\tau=\typ A_\ell$.}\prove{$X\cong L(Y)$ for some graph $Y$.}}
  \begin{proof}\pf\ $X$ is represented in $\typ A_\ell$ by Lemma~\ref{lem:4-transport} with $f$
   and \stepref{s414mR}; apply Theorem~\ref{thm:4.1}.~\qed\end{proof}
  \step{s414mD}{\case{$\tau=\typ D_\ell$.}\prove{$X\cong L(Y)$ for some graph $Y$, or $X\cong
   CP(b)$ for some $b$.}}
  \begin{proof}\pf\ $X$ is represented in $\typ D_\ell$ by Lemma~\ref{lem:4-transport} with $f$
   and \stepref{s414mR}, so it is a generalized line graph by Theorem~\ref{thm:4.2}, assertion 1;
   it is connected and regular by \asm[1]{0}, so apply
   Lemma~\ref{lem:4-regglg}.~\qed\end{proof}
  \step{s414mnotL}{$X\not\cong L(Y)$ for every graph $Y$.}
  \begin{proof}\pf\ Otherwise Lemma~\ref{lem:414-linegraph}, with \asm[1]{0} and
   \stepref{s414mdet}, gives $X\cong L(K_m)$ or $L(K_{m,m})$, against \asm[4]{0}.~\qed\end{proof}
  \qedstep
  \begin{proof}\pf\ In both cases \stepref{s414mA}, \stepref{s414mD}, the conclusion contradicts
   \stepref{s414mnotL} or \asm[4]{0}.~\qed\end{proof}
 \end{proof}
 \step{s414mE}{$\tau=\typ E_t$ with $t=\rho\in\{6,7,8\}$, and $(n,k,a,c)$ is one of
  $(10,3,0,1)$, $(15,8,4,4)$, $(16,10,6,6)$ (for $t=6$), $(16,6,2,2)$, $(21,10,5,4)$,
  $(27,16,10,8)$ (for $t=7$) and $(28,12,6,4)$ (for $t=8$).}
 \begin{proof}\pf\ By \stepref{s414mAD}, $\tau$ is one of $\typ E_6$, $\typ E_7$, $\typ E_8$, of
  ranks 6, 7, 8 (Lemma~\ref{lem:sec3-spans}), so $t=\rho$ by \stepref{s414mrk}. Of the 15
  tuples of \stepref{s414mCeighteen}, \stepref{s414mnotCP} removes six, and $\rho\in\{6,7,8\}$
  removes $(9,4,1,2,5)$ and $(10,6,3,4,5)$.~\qed\end{proof}
 \step{s414mint}{$\det(B_t)\,n/(k+2)$ is an integer, where $\det B_t$ is 3, 2, 1 for $t=6,7,8$.}
 \begin{proof}
  \step{s414mbasis}{\pick\ $\Zb_i=F^{-1}(\Zs_i)$, $1\le i\le t$, with the roots $\Zs_i$ of
   Computation C21; they lie in $\Phi(R)$ and form a basis of $W$ with Gram matrix $B_t$.}
  \begin{proof}\pf\ Computation~\ref{comp:C21} gives the $\Zs_i$, a basis of
   $\opk{span}\Phi(\typ E_t)$; $F^{-1}$ is a surjective linear isometry onto $W$ that maps
   $\Phi(\typ E_t)$ onto $\Phi(R)$ (\stepref{s414mtype} and
   \stepref{s414mspan}).~\qed\end{proof}
  \step{s414mwb}{$\Zw\in W$, $(\Zw,\Zb_i)\in\ints$ for every $i$, and $(\Zw,\Zw)=n/(k+2)$.}
  \begin{proof}\pf\ Lemma~\ref{lem:414-apex}, with \asm[1]{0} and \stepref{s414mdet}; the $\Zb_i$
   lie in $\Phi(R)\sbs \Lambda$ by \stepref{s414mbasis}.~\qed\end{proof}
  \qedstep
  \begin{proof}\pf\ Lemma~\ref{lem:414-dual}, with \stepref{s414mbasis} and \stepref{s414mwb};
   the determinants are those of Computation~\ref{comp:C21}.~\qed\end{proof}
 \end{proof}
 \step{s414mexcl}{$(n,k,a,c)$ is neither $(15,8,4,4)$ nor $(21,10,5,4)$.}
 \begin{proof}\pf\ These have $t=6$ and $t=7$ by \stepref{s414mE}, and
  $3\cdot\frac{15}{10}=\frac92$ and $2\cdot\frac{21}{12}=\frac72$ are not integers, against
  \stepref{s414mint}.~\qed\end{proof}
 \step{s414mclebsch}{\case{$(n,k,a,c)=(16,10,6,6)$.}\prove{$X$ is isomorphic to the Clebsch
  graph.}}
 \begin{proof}\pf\ Lemma~\ref{lem:414-SE2} with \asm[1]{0}.~\qed\end{proof}
 \step{s414mschlafli}{\case{$(n,k,a,c)=(27,16,10,8)$.}\prove{$X$ is isomorphic to the Schl\"afli
  graph.}}
 \begin{proof}\pf\ Lemma~\ref{lem:414-SE3} with \asm[1]{0}.~\qed\end{proof}
 \step{s414mpet}{\case{$(n,k,a,c)=(10,3,0,1)$.}\prove{$X$ is isomorphic to the Petersen graph.}}
 \begin{proof}
  \step{s414mpetsw}{$X\cong L(\Delta_6)^H$ for some $H\sbs\Ecal_6$.}
  \begin{proof}\pf\ $t=6$ by \stepref{s414mE}, $(\Zw,\Zw)=10/5=2$ by Lemma~\ref{lem:414-apex}, and
   $n=10=|\Ecal_6|$ by Computation~\ref{comp:414-Nc}; apply Lemma~\ref{lem:414-transversal} with
   \asm[1]{0} and $f$ of \stepref{s414mtype}.~\qed\end{proof}
  \qedstep
  \begin{proof}\pf\ $L(\Delta_6)^H\cong X$ is 3-regular by \stepref{s414mpetsw}; apply
   Lemma~\ref{lem:414-E6}.~\qed\end{proof}
 \end{proof}
 \step{s414mshr}{\case{$(n,k,a,c)=(16,6,2,2)$.}\prove{$X$ is isomorphic to the Shrikhande
  graph.}}
 \begin{proof}
  \step{s414mshrsw}{$X\cong L(\Delta_7)^H$ for some $H\sbs\Ecal_7$.}
  \begin{proof}\pf\ $t=7$ by \stepref{s414mE}, $(\Zw,\Zw)=16/8=2$ by Lemma~\ref{lem:414-apex},
   and $n=16=|\Ecal_7|$ by Computation~\ref{comp:414-Nc}; apply Lemma~\ref{lem:414-transversal}
   with \asm[1]{0} and $f$ of \stepref{s414mtype}.~\qed\end{proof}
  \qedstep
  \begin{proof}\pf\ $L(\Delta_7)^H\cong X$ is 6-regular by \stepref{s414mshrsw}, so it is the
   Shrikhande graph or $L(K_{4,4})$ by Lemma~\ref{lem:414-E7}; the latter is excluded by
   \asm[4]{0}.~\qed\end{proof}
 \end{proof}
 \step{s414mchang}{\case{$(n,k,a,c)=(28,12,6,4)$.}\prove{$X$ is isomorphic to one of the Chang
  graphs.}}
 \begin{proof}
  \step{s414mchangsw}{$X\cong L(\Delta_8)^H$ for some $H\sbs\Ecal_8$.}
  \begin{proof}\pf\ $t=8$ by \stepref{s414mE}, $(\Zw,\Zw)=28/14=2$ by Lemma~\ref{lem:414-apex},
   and $n=28=|\Ecal_8|$ by Computation~\ref{comp:414-Nc}; apply Lemma~\ref{lem:414-transversal}
   with \asm[1]{0} and $f$ of \stepref{s414mtype}.~\qed\end{proof}
  \qedstep
  \begin{proof}\pf\ $L(\Delta_8)^H\cong X$ is 12-regular by \stepref{s414mchangsw}, so it is $T(8)$ or
   a Chang graph by Lemma~\ref{lem:414-E8}; $T(8)=L(K_8)$ is excluded by
   \asm[4]{0}.~\qed\end{proof}
 \end{proof}
 \qedstep
 \begin{proof}\pf\ By \stepref{s414mE} and \stepref{s414mexcl}, the cases \stepref{s414mclebsch},
  \stepref{s414mschlafli}, \stepref{s414mpet}, \stepref{s414mshr} and \stepref{s414mchang} are
  exhaustive, and each gives one of the seven graphs; it is only one of them by
  Proposition~\ref{prop:414-seven}, assertion 3. Assertion 2 is
  Proposition~\ref{prop:414-seven}, assertions 1 and 2.~\qed\end{proof}
\end{proof}

\medskip\noindent\textbf{What the proof relies on.}
Besides Sections 2 to 4, the proof uses the computations C4 (with its transitivity on ordered
stars),
C9 (with its two certificate isomorphisms), C18, C20 (at the apex), C21 and C22, and the classification SE4 of
\path{cgss/proofs/reg8/reg8.tex}. All of it is in the Lean formalization
(\path{LineGraphs/Seidel/}, \path{LineGraphs/Computations/C18.lean}, \path{LineGraphs/Computations/C22.lean} and its data files): C18 as a
checked search, C22's fourteen isomorphism certificates with the clique profiles and counts, SE4
by a case split with checked certificates, and the restated remark as
\path{SimpleGraph.seidel_seven_classification} and \path{SimpleGraph.seidel_seven}.

\endgroup
